% ============================================================
%  Water Cycle and Resources — Experiment Booklet (one section per experiment)
%  Year 7 (Year 8 extension)  |  Victorian Curriculum F–10 Version 2.0
%  Compile with: xelatex three times, then python3 fill.py Water_Resources_Experiment_Booklet_Yr7
% ============================================================
\documentclass[11pt,a4paper]{article}
\newcommand{\unitname}{Water Cycle and Resources}\newcommand{\unitcaps}{WATER CYCLE AND RESOURCES}

% ── Fonts (compile with LuaLaTeX) ────────────────────────────────────────────
\usepackage{fontspec}
% Bundled fonts (in ./fonts/) so the document compiles anywhere with XeLaTeX/LuaLaTeX
\setmainfont{Lato}[                   % body — humanist sans (stands in for deck's Inter)
  Path=./fonts/, Extension=.ttf, Ligatures=TeX,
  UprightFont=Lato-Regular, ItalicFont=Lato-Italic,
  BoldFont=Lato-Bold, BoldItalicFont=Lato-BoldItalic ]
\newfontfamily\headingfont{Poppins}[  % display/headings — matches the slide decks
  Path=./fonts/, Extension=.ttf, Ligatures=TeX,
  UprightFont=Poppins-Regular, ItalicFont=Poppins-Italic,
  BoldFont=Poppins-Bold, BoldItalicFont=Poppins-BoldItalic ]
\setlength{\headheight}{22pt}

\usepackage[protrusion=true,expansion=false]{microtype}
\usepackage[a4paper, top=2cm, bottom=2cm, left=2cm, right=2cm]{geometry}
\usepackage{xcolor}
\usepackage{tikz}
\usetikzlibrary{arrows.meta, shapes.geometric, calc, positioning, decorations.pathreplacing}
\usepackage{tcolorbox}
\tcbuselibrary{skins, breakable, raster}
\usepackage{fancyhdr}
\usepackage{titlesec}
\usepackage{enumitem}
\usepackage{booktabs}
\usepackage{tabularx}
\usepackage{array}
\usepackage{colortbl}
\usepackage{multicol}
\usepackage{amsmath}
\usepackage{needspace}
% TikZ diagram macros are embedded below (just before \begin{document}) for easy editing.
\usepackage{longtable}
\usepackage{lastpage}
\usepackage{calc}
\usepackage{mdframed}
\usepackage{setspace}
\usepackage{ragged2e}
% PDF bookmarks: one per lesson (made in \sectiondiv) plus the cover and
% contents, and links from the contents table to each lesson.
% implicit=false: hyperref adds no anchors of its own (an anchor on each
% list item nudged lists down 1.5 mm, which would undo the page fill).
% hidelinks keeps the links invisible, so the printed pages are unchanged.
\usepackage[implicit=false, hidelinks,
  pdftitle={Water Cycle and Resources — Experiment Booklet},
  pdfsubject={Years 7–8 Science, Victorian Curriculum F–10 v2.0}]{hyperref}
\usepackage{bookmark} % writes the outline itself, which implicit=false switches off in hyperref
% hyperref 7.01 with implicit=false leaves \Hy@PageAnchorSlide undefined, which
% throws at \end{document} and stops fill.py (it builds with -halt-on-error).
\makeatletter\providecommand\Hy@PageAnchorSlide{}\makeatother

% ── Colour Palette ───────────────────────────────────────────────────────────
% Remapped to the slide-deck palette (names kept so body + diagrams cascade)
\definecolor{NavyBlue}   {HTML}{16305C}
\definecolor{MidBlue}    {HTML}{2E6FE0}
\definecolor{LightBlue}  {HTML}{E8F1FF}
\definecolor{TealGreen}  {HTML}{0F6E63}
\definecolor{LightGreen} {HTML}{E0F7F3}
\definecolor{GoldDark}   {HTML}{C77F12}
\definecolor{GoldLight}  {HTML}{FFF3DC}
\definecolor{PurpleDark} {HTML}{5B3FC0}
\definecolor{PurpleLight}{HTML}{EFE9FF}
\definecolor{OrangeDark} {HTML}{C9483D}
\definecolor{OrangeLight}{HTML}{FFE9E6}
\definecolor{MidGrey}    {HTML}{AEBED4}
\definecolor{LightGrey}  {HTML}{F2F5FA}
\definecolor{DarkGrey}   {HTML}{3A4759}
\definecolor{TableHead}  {HTML}{16305C}
\definecolor{RowAlt}     {HTML}{EEF4FF}
% Bright accent tints (deck) for bars, rules and chips
\definecolor{Blue}   {HTML}{2E6FE0}
\definecolor{Teal}   {HTML}{18A999}
\definecolor{Amber}  {HTML}{F5A623}
\definecolor{Coral}  {HTML}{FF6B5E}
\definecolor{Violet} {HTML}{7C5CE0}
\definecolor{Sky}    {HTML}{4DA3FF}
\definecolor{Ink}    {HTML}{16223A}
\definecolor{Line}   {HTML}{DDE5F0}
\definecolor{Paper}  {HTML}{FBFCFE}

% ── Page Style ───────────────────────────────────────────────────────────────
\pagestyle{fancy}
\fancyhf{}
\renewcommand{\headrulewidth}{0.4pt}
\renewcommand{\headrule}{\hbox to\headwidth{\color{NavyBlue}\leaders\hrule height \headrulewidth\hfill}}
\fancyhead[L]{\headingfont\small\color{NavyBlue}\unitname{} \textbar{} Experiment Booklet}
\fancyhead[R]{\headingfont\small\color{DarkGrey}Year 7 \textbar{} Victorian Curriculum F--10 v2.0}
\fancyfoot[C]{\headingfont\small\color{DarkGrey}\thepage\ \textbar{} \pageref{LastPage}}
\fancypagestyle{plain}{%
  \fancyhf{}
  \fancyfoot[C]{\headingfont\small\color{DarkGrey}\thepage\ \textbar{} \pageref{LastPage}}
  \renewcommand{\headrulewidth}{0pt}
}

% ── Section Formatting ───────────────────────────────────────────────────────
\titleformat{\section}
  {\headingfont\color{NavyBlue}\Large\bfseries}
  {}{0em}{}
  [{\color{Teal}\rule[0.6ex]{2.2em}{3pt}}{\color{Line}\rule[0.6ex]{\dimexpr\linewidth-2.2em\relax}{1pt}}]
\titlespacing{\section}{0pt}{18pt}{8pt}

\titleformat{\subsection}
  {\headingfont\color{MidBlue}\large\bfseries}
  {\thesubsection}{0.6em}{}
\titlespacing{\subsection}{0pt}{12pt}{4pt}

\titleformat{\subsubsection}
  {\headingfont\color{NavyBlue}\normalsize\bfseries}
  {}{0em}{}
\titlespacing{\subsubsection}{0pt}{8pt}{2pt}

% Deck-style eyebrow / kicker (optional use)
\newcommand{\kicker}[2][Teal]{%
  {\headingfont\bfseries\footnotesize\color{#1}%
   \textcolor{#1}{\rule[0.18ex]{1.6em}{2.2pt}}\hspace{0.5em}\MakeUppercase{#2}\par}%
  \vspace{2pt}}

% ── Custom tcolorbox Environments ────────────────────────────────────────────

% Section divider banner
\newtcolorbox{sectionbanner}[2]{
  enhanced, breakable=false,
  colback=NavyBlue, colframe=NavyBlue,
  fontupper=\color{white}\bfseries\large,
  left=4mm, right=4mm, top=3mm, bottom=3mm,
  title={\color{white}\Huge\bfseries #1},
  attach boxed title to top left={yshift=-2mm, xshift=4mm},
  boxed title style={colback=NavyBlue, colframe=NavyBlue},
  after upper={\par\smallskip\normalfont\normalsize\color{white!80!NavyBlue}\textit{#2}}
}

% Key info box (blue)
\newtcolorbox{infobox}[1]{
  breakable, arc=2.5mm, boxrule=0.5pt,
  colback=LightBlue, colframe=Blue!40,
  fonttitle=\headingfont\bfseries, coltitle=MidBlue,
  title={#1},
  left=4mm, right=3.5mm, top=2.5mm, bottom=2.5mm
}

% Scaffold box (yellow)
\newtcolorbox{scaffold}[1]{
  breakable, arc=2.5mm, boxrule=0.5pt,
  colback=GoldLight, colframe=Amber!40,
  fonttitle=\headingfont\bfseries, coltitle=GoldDark,
  title={#1},
  left=4mm, right=3.5mm, top=2.5mm, bottom=2.5mm
}

% Worked example box
\newtcolorbox{workedexbox}[1]{
  breakable, arc=2.5mm, boxrule=0.5pt,
  colback=white, colframe=Amber!50,
  fonttitle=\headingfont\bfseries, coltitle=GoldDark,
  title={#1},
  left=4mm, right=3.5mm, top=2.5mm, bottom=2.5mm
}
% Worked examples: \begin{workedex}{title} -> "Worked Example N.M --- title",
% auto-numbered per section and labelled we:N.M. Never type the number.
\newcounter{wenum}[section]
\renewcommand{\thewenum}{\thesection.\arabic{wenum}}
\newenvironment{workedex}[1]
  {\refstepcounter{wenum}\label{we:\thewenum}\begin{workedexbox}{Worked Example \thewenum\ --- #1}}
  {\end{workedexbox}}

% We do: straight after each worked example, a near-identical question the
% class works through together, with its own writing space. Numbered to
% match the worked example above it and labelled wd:N.M.
%   \begin{wedo}{title} question ... \answerbox{h} \end{wedo}
\newtcolorbox{wedobox}[1]{
  enhanced, breakable=false, arc=2.5mm, boxrule=0.5pt,
  colback=white, colframe=TealGreen!55,
  fonttitle=\headingfont\bfseries, coltitle=TealGreen!75!black,
  title={#1},
  left=4mm, right=3.5mm, top=2.5mm, bottom=2.5mm
}
\newenvironment{wedo}[1]
  {\par\needspace{9\baselineskip}\begin{wedobox}{We do \thewenum\ --- #1}\label{wd:\thewenum}\ignorespaces}
  {\end{wedobox}}

% Year 7 box (blue)
\newtcolorbox{yr7box}{
  breakable, arc=2.5mm, boxrule=0.5pt,
  colback=LightBlue, colframe=Blue!40,
  left=4mm, right=3.5mm, top=2.5mm, bottom=2.5mm
}

% Year 8 extension box (violet)
\newtcolorbox{yr8box}{
  breakable, arc=2.5mm, boxrule=0.5pt,
  colback=PurpleLight, colframe=Violet!40,
  left=4mm, right=3.5mm, top=2.5mm, bottom=2.5mm
}

% Answer lines box
\newcommand{\answerbox}[1]{%
  \begin{tcolorbox}[
    enhanced, breakable, arc=2mm,
    colback=Paper, colframe=Line,
    left=2mm, right=2mm, top=1mm, bottom=1mm,
    boxrule=1pt,
    height=#1
  ]
  \end{tcolorbox}%
}

% Experiment header
\newtcolorbox{expheader}[3]{
  breakable=false, arc=2.5mm,
  colback=RowAlt, colframe=Amber!50, boxrule=0.5pt,
  fonttitle=\headingfont\bfseries\large, coltitle=NavyBlue,
  title={Experiment #1: #2},
  left=4mm, right=3.5mm, top=2.5mm, bottom=2.5mm,
  after upper={\par\smallskip
    {\headingfont\bfseries\color{GoldDark}Aim:}\ \textit{#3}},
}

% Reading comprehension passage box (green)
\newtcolorbox{readingbox}[1]{
  breakable, arc=2.5mm, boxrule=0.5pt,
  colback=white, colframe=Teal!50,
  fonttitle=\headingfont\bfseries, coltitle=TealGreen,
  title={#1},
  left=4mm, right=3.5mm, top=2.5mm, bottom=2.5mm,
  fontupper=\normalsize
}

% Skill-builder box (for beginner scaffolding)
\newtcolorbox{skillbox}[1]{
  breakable, arc=2.5mm, boxrule=0.5pt,
  colback=RowAlt, colframe=Blue!40,
  fonttitle=\headingfont\bfseries, coltitle=MidBlue,
  title={#1},
  left=4mm, right=3.5mm, top=2.5mm, bottom=2.5mm
}

% ── Custom Commands ───────────────────────────────────────────────────────────

% Question: \question{marks}{text}  (marks can be empty)
% Numbered automatically as Q<section>.<n>, reset every section.
% Do NOT write numbers by hand.
\newcounter{qnum}[section]
\renewcommand{\theqnum}{\thesection.\arabic{qnum}}
\newcommand{\question}[2]{%
  \refstepcounter{qnum}%
  \needspace{3\baselineskip}%
  \medskip
  \noindent\label{q:\theqnum}\textbf{Q\theqnum}\ifx&#1&\else\ \textcolor{DarkGrey}{(#1 mark\ifnum#1=1\else s\fi)}\fi\quad #2%
  \medskip
}

% Inline marks note
\newcommand{\markscount}[1]{\hfill\textcolor{DarkGrey}{\textit{(#1 mark\ifnum#1=1\else s\fi)}}}

% ── Practice set counter and question macros ─────────────────────────────────
\newcounter{pqnum}
\newcommand{\practiceset}[1]{%
  \needspace{6\baselineskip}
  \setcounter{pqnum}{0}%
  \vspace{6pt}
  {\headingfont\bfseries\large\color{GoldDark} Practice Set --- #1}\\[-2pt]
  {\color{Amber}\rule{\linewidth}{1pt}}\\[2pt]
}
% \pq{question text}{working box height}
\newcommand{\pq}[2]{%
  \stepcounter{pqnum}%
  \needspace{4\baselineskip}
  \vspace{5pt}
  \noindent\textbf{\thepqnum.}\quad #1\par
  \vspace{2pt}
  \answerbox{#2}
}
% Two short questions side by side, each with its own box
\newcommand{\pqpair}[4]{%
  \needspace{4\baselineskip}
  \vspace{5pt}
  \noindent
  \begin{minipage}[t]{0.48\linewidth}
    \stepcounter{pqnum}\textbf{\thepqnum.}\quad #1\par\vspace{2pt}
    \begin{tcolorbox}[enhanced,arc=2mm,colback=Paper,colframe=Line,left=2mm,right=2mm,
      top=1mm,bottom=1mm,boxrule=1pt,height=#2]\end{tcolorbox}
  \end{minipage}\hfill
  \begin{minipage}[t]{0.48\linewidth}
    \stepcounter{pqnum}\textbf{\thepqnum.}\quad #3\par\vspace{2pt}
    \begin{tcolorbox}[enhanced,arc=2mm,colback=Paper,colframe=Line,left=2mm,right=2mm,
      top=1mm,bottom=1mm,boxrule=1pt,height=#4]\end{tcolorbox}
  \end{minipage}
}

% ── Comprehension question macro (compact, lined) ────────────────────────────
\newcounter{rcnum}
\newcommand{\rcreset}{\setcounter{rcnum}{0}}
\newcommand{\rcq}[2]{%
  \stepcounter{rcnum}%
  \needspace{3\baselineskip}
  \vspace{4pt}
  \noindent\textbf{C\thercnum.}\quad #1\par
  \vspace{2pt}
  \linedbox{#2}
}

% Section divider (big numbered banner): \sectiondiv{title}{summary}
%
% This IS the section: it steps the real section counter, so subsections,
% tasks, questions and figures inside it are numbered N.M automatically.
% Put \label{sec:N} straight after it. Never hand-type a section number.
\newcommand{\sectiondiv}[2]{%
  \needspace{5\baselineskip}%
  \refstepcounter{section}%
  % PDF bookmark for the lesson; its anchor, lesson.N.1, is what the contents table links to
  \pdfbookmark[1]{Lesson \thesection{} — #1}{lesson.\thesection}%
  \begin{tcolorbox}[
    enhanced, breakable=false, arc=3.5mm,
    colback=NavyBlue, colframe=NavyBlue, boxrule=0pt,
    sidebyside, sidebyside align=center,
    lefthand width=2.4cm,
    left=2mm, right=4mm, top=4mm, bottom=4mm
  ]
    \begin{center}
      {\headingfont\bfseries\color{white!35!NavyBlue}\fontsize{50}{50}\selectfont \thesection}
    \end{center}
    \tcblower
    {\headingfont\bfseries\color{white}\LARGE #1}\\[3pt]
    {\itshape\color{white!75!NavyBlue}\small #2}
  \end{tcolorbox}%
}

% Tasks: \task, \task[a] (part a of a NEW task), \taskpart{b} (next part
% of the same task), \exttask (Year 8 "Extension" — same counter).
% Numbered Task <section>.<n>, reset every section. Do NOT type numbers.
% Every task, question and figure also gets an automatic label
% (task:1.7b, q:2.4, fig:1.2, fig:E2.1) so page refs come from the .aux.
\newcounter{tasknum}[section]
\renewcommand{\thetasknum}{\thesection.\arabic{tasknum}}
\NewDocumentCommand{\task}{o}{\refstepcounter{tasknum}\label{task:\thetasknum\IfValueT{#1}{#1}}\textbf{Task \thetasknum\IfValueT{#1}{#1}}}
\newcommand{\taskpart}[1]{\label{task:\thetasknum#1}\textbf{Task \thetasknum#1}}
\newcommand{\exttask}{\refstepcounter{tasknum}\label{task:\thetasknum}\textbf{Extension \thetasknum}}

% Figure captions: \figcap[fig:label]{caption} -> "Figure N.M --- caption"
% \expfigcap{caption} -> "Figure EN.M" (experiment pages). Both auto-numbered.
\newcounter{fignum}[section]
\renewcommand{\thefignum}{\thesection.\arabic{fignum}}
\newcounter{expfignum}[section]
\renewcommand{\theexpfignum}{E\thesection.\arabic{expfignum}}
\newcommand{\figcap}[2][]{%
  \begin{center}
    \refstepcounter{fignum}\label{fig:\thefignum}\if\relax\detokenize{#1}\relax\else\label{#1}\fi
    \small\itshape\color{DarkGrey} Figure \thefignum\ --- #2
  \end{center}%
}
\newcommand{\expfigcap}[1]{%
  \begin{center}
    \refstepcounter{expfignum}\label{fig:\theexpfignum}%
    \small\itshape\color{DarkGrey} Figure \theexpfignum\ --- #1
  \end{center}%
}

% Diagram include
\newcommand{\diagramfig}[3]{%  {filename}{width fraction}{caption}
  \begin{center}
    \includegraphics[width=#2\linewidth]{#1}\\[4pt]
    \small\itshape\color{DarkGrey} #3
  \end{center}%
}

% Table column types
\newcolumntype{L}[1]{>{\RaggedRight\arraybackslash}p{#1}}
\newcolumntype{C}[1]{>{\Centering\arraybackslash}p{#1}}

% Table header colour
\newcommand{\thead}[1]{\cellcolor{TableHead}\color{white}\bfseries #1}
\newcommand{\talt}{\rowcolor{RowAlt}}
\newcommand{\tmid}{\rowcolor{MidGrey}}

% Inline year tags
\newcommand{\tagyr}[1]{%
  \tikz[baseline=-0.5ex]\node[
    rounded corners=2pt,
    fill=#1!20, draw=#1, text=#1,
    inner xsep=4pt, inner ysep=1pt,
    font=\scriptsize\bfseries
  ]{Yr #1 only};}

% Draw answer lines
\newcommand{\anslines}[1]{%
  \vspace{2pt}%
  \foreach \n in {1,...,#1}{%
    \par\noindent\rule{\linewidth}{0.3pt}\vspace{3pt}%
  }%
}

% ── Document ─────────────────────────────────────────────────────────────────
% ── Optional-lesson note, shown under the lesson banner ──────────────────────
\newcommand{\optionalnote}[1]{%
  % the title is a bold lead-in on the first line rather than a title bar,
  % which keeps Lesson 11's reading on the lesson's first page
  \begin{tcolorbox}[enhanced, breakable=false, arc=2.5mm, boxrule=0.5pt,
    colback=PurpleLight, colframe=Violet!50,
    left=4mm, right=3.5mm, top=2mm, bottom=2mm]
  {\headingfont\bfseries\color{PurpleDark}Optional lesson \textbullet{} Year 8 extension.}\enspace #1
  \end{tcolorbox}\vspace{2pt}}


% ══════════════════════════════════════════════════════════════════════════════
%  LESSON STRUCTURE (lesson booklet, Oct 2026)
%  Every lesson runs in the same four parts, in this order:
%    \lessonpart{A}{Reading and comprehension}{Teal}   \label{read:N}
%    \lessonpart{B}{Theory and examples}{Blue}         \label{theory:N}
%    \lessonpart{C}{Practice}{Amber}                    \label{prac:N}
%    \lessonpart{D}{Above level}{Violet}                \label{ext:N}
%  Above-level / higher-order questions live ONLY in Part D, at the end.
% ══════════════════════════════════════════════════════════════════════════════
\newcommand{\lessonpart}[3]{%
  \par\if#1A\needspace{10\baselineskip}\vspace{10pt}\else
    \if#1C\notesfill\fi\clearpage\fi
  \begin{tcolorbox}[enhanced, breakable=false, arc=2mm, boxrule=0pt,
    colback=#3!14, borderline west={4pt}{0pt}{#3},
    left=4mm, right=3mm, top=1.6mm, bottom=1.6mm]
  {\headingfont\bfseries\color{#3!60!black}\large Part #1 \textbullet{} #2}
  \end{tcolorbox}\vspace{2pt}}

% Key vocabulary at the start of each lesson: students copy each definition
% from the vocabulary slide in the lesson's starter. First row = red scaffold.
\newcommand{\vocabhead}{%
  \par\vspace{6pt}%
  \begin{tcolorbox}[enhanced, breakable=false, arc=2mm, boxrule=0pt,
    colback=NavyBlue!10, borderline west={4pt}{0pt}{NavyBlue},
    left=4mm, right=3mm, top=1.6mm, bottom=1.6mm]
  {\headingfont\bfseries\color{NavyBlue}\large Key vocabulary}
  \end{tcolorbox}\vspace{2pt}%
  \noindent Copy the definition of each word from the slide. The first one has been done for you.\par}

% Curriculum codes under a lesson banner
\newcommand{\lessoncodes}[1]{%
  {\small\headingfont\color{DarkGrey}Victorian Curriculum 2.0:\ \textbf{#1}}\par\vspace{4pt}}

% Two short tasks side by side, each with its own working box
% \taskpair{text}{box height}{text}{box height}
% The two boxes line up: the shorter question is padded by the difference in
% text height (measured in a scratch box, task counter restored, no label).
\newsavebox\tpbox
\newlength\tpA \newlength\tpB
\newcount\tpsaved
\newcommand\tpmeasure[2]{%
  \tpsaved=\value{tasknum}%
  \sbox\tpbox{\vbox{\hsize=0.485\linewidth \linewidth=\hsize \parskip=0pt
    \renewcommand\label[1]{}\task #1\par}}%
  \setlength{#2}{\dimexpr\ht\tpbox+\dp\tpbox\relax}%
  \setcounter{tasknum}{\tpsaved}}
\newcommand{\taskpair}[4]{%
  \par\needspace{9\baselineskip}\vspace{4pt}%
  \tpmeasure{#1}{\tpA}\tpmeasure{#3}{\tpB}%
  \noindent
  \edef\abgrp{\number\numexpr\value{abox}+1\relax}%
  \begin{minipage}[t]{0.485\linewidth}
    \task #1\par\ifdim\tpB>\tpA\vspace{\dimexpr\tpB-\tpA\relax}\fi\vspace{2pt}\answerbox{#2}
  \end{minipage}\hfill
  \begin{minipage}[t]{0.485\linewidth}
    \task #3\par\ifdim\tpA>\tpB\vspace{\dimexpr\tpA-\tpB\relax}\fi\vspace{2pt}\answerbox{#4}
  \end{minipage}\par\gdef\abgrp{0}\vspace{4pt}}
% Two parts of one question side by side, with their answer boxes level:
% \ptpair{a}{part (a) text}{box height}{b}{part (b) text}{box height}.
% The shorter part's text gets space below it to match the longer one.
\newcommand\ptmeasure[3]{%
  \sbox\tpbox{\vbox{\hsize=0.485\linewidth \linewidth=\hsize \parskip=0pt \pt{#1}#2\par}}%
  \setlength{#3}{\dimexpr\ht\tpbox+\dp\tpbox\relax}}
\newcommand{\ptpair}[6]{%
  \ptmeasure{#1}{#2}{\tpA}\ptmeasure{#4}{#5}{\tpB}%
  \par\noindent
  \edef\abgrp{\number\numexpr\value{abox}+1\relax}%
  \begin{minipage}[t]{0.485\linewidth}
    \pt{#1}#2\par\ifdim\tpB>\tpA\vspace{\dimexpr\tpB-\tpA\relax}\fi\answerbox{#3}
  \end{minipage}\hfill
  \begin{minipage}[t]{0.485\linewidth}
    \pt{#4}#5\par\ifdim\tpA>\tpB\vspace{\dimexpr\tpA-\tpB\relax}\fi\answerbox{#6}
  \end{minipage}\par\gdef\abgrp{0}}

% Answer boxes never split across a page (a half box on its own page is useless)
% ── Page filling ────────────────────────────────────────────────────────────
% Every writing box gets an id. Its page is logged to \jobname.abp; fill.py
% measures the empty space left at the foot of each page and shares it out
% among that page's boxes via \jobname.abx (\abset{id}{extra}). Side-by-side
% boxes share a group so they grow together.
\newcounter{abox}
\newwrite\abpos \immediate\openout\abpos=\jobname.abp
\newcommand{\abset}[2]{\expandafter\gdef\csname abx@#1\endcsname{#2}}
\InputIfFileExists{\jobname.abx}{}{}
\def\abgrp{0}
\newcommand{\abnext}{\stepcounter{abox}%
  \edef\abcur{\ifcsname abx@\theabox\endcsname\csname abx@\theabox\endcsname\else 0pt\fi}}
\newcommand{\ablog}{\edef\abtmp{\noexpand\write\abpos{\theabox\space\abgrp\space\noexpand\thepage}}\abtmp}
\newcommand{\abh}[1]{\dimexpr#1+\abcur\relax}
\renewcommand{\answerbox}[1]{%
  \abnext
  \begin{tcolorbox}[enhanced, breakable=false, arc=2mm, colback=Paper, colframe=Line,
    left=2mm, right=2mm, top=1mm, bottom=1mm, boxrule=1pt, height=\abh{#1}]
  \end{tcolorbox}\ablog%
}
% \linedbox{height}: an answer box ruled in dotted thirds. Lines are \thirdrow/3 apart and run dotted, dotted,
% solid (the baseline), repeating, and end on one dotted line. The number of
% rows is worked out from the box's final height, so the ruling still fills the
% box after fill.py has grown it; the rows are centred in the box.
\newlength\thirdrow \setlength\thirdrow{9mm}
\tikzset{thirdsolid/.style={draw=DarkGrey!70, line width=0.5pt},
         thirddot/.style={draw=DarkGrey!70, line width=0.6pt, dash pattern=on 0.9pt off 1.6pt}}
\newlength\lbH \newlength\lbO \newcount\lbN
\newcommand{\linedbox}[1]{%
  \abnext
  % interior height = box height less the 1pt frame top and bottom
  \setlength\lbH{\dimexpr\abh{#1}-2pt\relax}%
  % whole rows that fit with at least 3 mm spare top and bottom (at least one)
  \lbN=\numexpr(\dimexpr\lbH-6mm-\thirdrow/2\relax)/\thirdrow\relax
  \ifnum\lbN<1 \lbN=1 \fi
  \setlength\lbO{\dimexpr(\lbH-\lbN\thirdrow)/2\relax}%
  \begin{tcolorbox}[enhanced, breakable=false, arc=2mm, colback=Paper, colframe=Line,
    left=2mm, right=2mm, top=1mm, bottom=1mm, boxrule=1pt, height=\abh{#1},
    underlay={\begin{tcbclipinterior}
      \foreach \i in {1,...,\the\lbN}{
        \draw[thirddot] ([xshift=3mm,yshift=-\lbO-(\i-1)*\thirdrow]interior.north west) -- ([xshift=-3mm,yshift=-\lbO-(\i-1)*\thirdrow]interior.north east);
        \draw[thirddot] ([xshift=3mm,yshift=-\lbO-(\i-1)*\thirdrow-\thirdrow/3]interior.north west) -- ([xshift=-3mm,yshift=-\lbO-(\i-1)*\thirdrow-\thirdrow/3]interior.north east);
        \draw[thirdsolid] ([xshift=3mm,yshift=-\lbO-(\i-1)*\thirdrow-2*\thirdrow/3]interior.north west) -- ([xshift=-3mm,yshift=-\lbO-(\i-1)*\thirdrow-2*\thirdrow/3]interior.north east);}
      \draw[thirddot] ([xshift=3mm,yshift=-\lbO-\lbN\thirdrow]interior.north west) -- ([xshift=-3mm,yshift=-\lbO-\lbN\thirdrow]interior.north east);
    \end{tcbclipinterior}}]
  \end{tcolorbox}\ablog%
}
% Every written-response box is ruled the same way (Part A comprehension, We do,
% Part C tasks and Part D extensions, including side-by-side pairs). Space for
% mathematical working uses \workbox{height} instead and stays plain, as do
% \drawbox and the write-in tables.
\let\workbox\answerbox   % plain box: mathematical working (calculations) stays unruled
\let\answerbox\linedbox
% \drawbox{height}{contents}: white box for a diagram or drawing
\newcommand{\drawbox}[2]{%
  \abnext
  \begin{tcolorbox}[colback=white, colframe=MidGrey, boxrule=0.5pt, height=\abh{#1}]#2\end{tcolorbox}\ablog}
% Notes box at the end of Part B: fills whatever is left of the page (if worth it)
\newcommand{\notesfill}{\par
  \ifdim\dimexpr\pagegoal-\pagetotal\relax>3.2cm
    \vspace*{\fill}%
    \begin{tcolorbox}[enhanced, breakable=false, arc=2mm, colback=Paper, colframe=Line,
      left=2mm, right=2mm, top=1mm, bottom=1mm, boxrule=1pt,
      height=\dimexpr\pagegoal-\pagetotal-1.0cm\relax]
      {\small\headingfont\color{DarkGrey}Notes}
    \end{tcolorbox}%
  \fi}



% Centred table: wrap every content table in \begin{ctab} ... \end{ctab}
% (figures already centre themselves inside their \diag... macros)
% ctab also registers write-in tables (S column) for page filling: the extra
% height is per row, logged with the row count.
\newenvironment{ctab}{\par\smallskip\begingroup\centering
  \abnext\let\tbcur\abcur\global\wtrows0 }{\par
  \ifnum\wtrows>0 \edef\abtmp{\noexpand\write\abpos{\theabox\space\abgrp\space\noexpand\thepage\space\the\wtrows}}\abtmp\fi
  \endgroup\smallskip}

% ── Booklet formatting rules (Oct 2026) ─────────────────────────────────────
% 1. No marks are shown anywhere in the booklet.
\renewcommand{\markscount}[1]{}
% 2. A task's heading sits on its own line; its text starts on the next
%    line, left aligned. Parts (a), (b)… each start a new line: use \pt{a}.
\RenewDocumentCommand{\task}{o}{\par\noindent\refstepcounter{tasknum}\label{task:\thetasknum\IfValueT{#1}{#1}}\textbf{Task \thetasknum\IfValueT{#1}{#1}}\par\nobreak\noindent\ignorespaces}
\renewcommand{\taskpart}[1]{\par\noindent\label{task:\thetasknum#1}\textbf{Task \thetasknum#1}\par\nobreak\noindent\ignorespaces}
\renewcommand{\exttask}{\par\noindent\refstepcounter{tasknum}\label{task:\thetasknum}\textbf{Extension \thetasknum}\par\nobreak\noindent\ignorespaces}
\newcommand{\pt}[1]{\par\noindent\textbf{(#1)}\ \ignorespaces}
% Workstation pages (Experiment 1): numbered method steps with a hanging
% indent, and a one-line safety reminder under the materials list.
\newcommand{\wsstep}[1]{\par\noindent\hangindent=1.6em\hangafter=1\makebox[1.6em][l]{\textbf{#1}}\ignorespaces}
\newcommand{\wssafety}[1]{\par\smallskip{\small\textcolor{ScaffoldRed}{\headingfont\bfseries Safety:}\ #1}\par}
% 3. Tables have vertical rules between columns and down both sides
%    (write | between and around the columns, e.g. {|L{3cm}|C{2cm}|}).
%    Keep the normal padding at the outer edges (no @{} there): colortbl paints
%    header colour one \tabcolsep past each edge, so trimming the padding makes
%    the header overhang the grid.
\arrayrulecolor{MidGrey}
\setlength{\arrayrulewidth}{0.6pt}
% 4. Write-in tables: W/Z are left/centred columns that centre text
%    vertically; put the S column first (see below).
\newcolumntype{W}[1]{>{\RaggedRight\arraybackslash}m{#1}}
\newcolumntype{Z}[1]{>{\Centering\arraybackslash}m{#1}}
\def\tbcur{0pt}
\newcount\wtrows
\newcommand{\wh}{\global\advance\wtrows1 \rule[\dimexpr-0.75cm-\tbcur/2\relax]{0pt}{\dimexpr1.5cm+\tbcur\relax}}
% S = zero-width first column that makes every row of a write-in table
% at least 1.5 cm tall with its text vertically centred. Start each data
% row with "&"; start the header row with "\hd &" so it stays slim.
\newif\ifnowh
\newcolumntype{S}{@{}>{}c<{\ifnowh\else\wh\fi\global\nowhfalse}@{\hspace{\tabcolsep}}}
\newcommand{\hd}{\global\nowhtrue}
% 5. The first row of a multi-row write-in table is a worked scaffold in red.
\definecolor{ScaffoldRed}{HTML}{C0392B}
\newcommand{\scaf}[1]{\textcolor{ScaffoldRed}{\itshape #1}}
% 6. Tables are a light grid (rules between rows and columns) on white —
%    row shading is off, because colortbl paints shading over column rules.
\renewcommand{\talt}{}
\renewcommand{\tmid}{}
% ══════════════════════════════════════════════════════════════════════════════
%  EXPERIMENT BOOKLET LAYER (Nov 2026)
%  From 2027 every Year 7 unit collects its experiments into one A4 experiment
%  booklet that puts the experimental techniques and skills first. Every skill
%  comes from the faculty Science Skills continuum (Years 7--8 column) and is
%  named the same way everywhere: a code (QP1, PC4, PA2 …) on a chip coloured by
%  its sub-strand. Copy this block into the next unit's experiment booklet.
%
%  Each experiment runs in the same order:
%    \experimentdiv{title}{summary}{lesson}\label{sec:N}
%    key vocabulary · learning intentions · \begin{skillsfocus}
%    \exppart*{Technique spotlight}{Blue}   (technique, worked example, We do)
%    \exppart{Plan and safety}{Teal}        (aim, prediction, variables, risks)
%    \exppart{Method}{Teal}
%    \exppart{Results}{Amber}
%    \exppart{Analysis and evaluation}{Coral}
%    \exppart{Skill check}{NavyBlue}
% ══════════════════════════════════════════════════════════════════════════════
\usetikzlibrary{arrows.meta,calc,positioning}
\usepackage{pifont}

\newenvironment{keepext}{\par\begin{tcolorbox}[blankest, breakable=false]}{\end{tcolorbox}\par}
% graph-paper box: \graphgrid{height}
\newcommand{\graphgrid}[1]{%
  \begin{tcolorbox}[enhanced, breakable=false, colback=white, colframe=MidGrey, boxrule=0.5pt, height=#1,
    underlay={\begin{tcbclipinterior}\draw[step=5mm, Line, very thin] (interior.south west) grid (interior.north east);
      \draw[step=1cm, MidGrey!70, thin] (interior.south west) grid (interior.north east);\end{tcbclipinterior}}]
  \end{tcolorbox}}

% ── Skills (faculty Science Skills continuum, Years 7–8) ─────────────────────
% \defskill{code}{name}{colour}{VC code}{what it means}
\makeatletter
\newcommand{\defskill}[5]{%
  \expandafter\gdef\csname sk@n@#1\endcsname{#2}%
  \expandafter\gdef\csname sk@c@#1\endcsname{#3}%
  \expandafter\gdef\csname sk@v@#1\endcsname{#4}%
  \expandafter\gdef\csname sk@d@#1\endcsname{#5}}
\newcommand{\skname}[1]{\csname sk@n@#1\endcsname}
\newcommand{\skcode}[1]{\csname sk@v@#1\endcsname}
\newcommand{\skdesc}[1]{\csname sk@d@#1\endcsname}
% \sk{PC4}: the chip — code and name on the sub-strand colour
\newcommand{\sk}[1]{%
  \edef\sk@col{\csname sk@c@#1\endcsname}%
  \tikz[baseline=(c.base)]\node[rounded corners=2.5pt, fill=\sk@col!16, draw=\sk@col, line width=0.6pt,
    text=\sk@col!55!black, inner xsep=4pt, inner ysep=2pt, font=\headingfont\bfseries\scriptsize]
    (c) {#1\,\textperiodcentered\,\skname{#1}};}
% \skr{PC4}: the chip flush right at the end of a paragraph; drops to its own line if it does not fit
\newcommand{\skr}[1]{\unskip\nobreak\hfill\penalty50\hskip1em\null\nobreak\hfill\sk{#1}}
% \skc{PC4}: the code alone, for tight table cells
\newcommand{\skc}[1]{%
  \edef\sk@col{\csname sk@c@#1\endcsname}%
  \tikz[baseline=(c.base)]\node[rounded corners=2.5pt, fill=\sk@col!16, draw=\sk@col, line width=0.6pt,
    text=\sk@col!55!black, inner xsep=3pt, inner ysep=2pt, font=\headingfont\bfseries\scriptsize]
    (c) {#1};}
\makeatother

% Questioning and predicting — VC2S8I01
\defskill{QP1}{Testable hypothesis}{Teal}{VC2S8I01}{A prediction an experiment can test, naming both variables and the direction of change: \emph{If the [IV] is \dots, the [DV] will \dots}; a reason is optional}
\defskill{QP2}{Formulate an aim}{Teal}{VC2S8I01}{One sentence naming both variables: \emph{To investigate the effect of changing [IV] on [DV]}}
% Planning and conducting — VC2S8I02, VC2S8I03
\defskill{PC1}{Plan and follow a method}{Blue}{VC2S8I02}{Follow, and later write, numbered steps in order: named equipment, quantities with units, no personal pronouns and a labelled diagram, so that someone else could repeat it (repeatable)}
\defskill{PC2}{Identify variables}{Blue}{VC2S8I02}{Name the independent variable (what you change), the dependent variable (what you measure) and at least two controlled variables (what you keep the same)}
\defskill{PC3}{Risks and safety}{Blue}{VC2S8I02}{Name each risk, how to prevent it and how to respond if it happens; follow the lab rules}
\defskill{PC4}{Precision}{Blue}{VC2S8I02}{Choose equipment with a fine enough scale, read it to its smallest division, and repeat measurements to see how closely they agree}
\defskill{PC5}{Ethics}{Blue}{VC2S8I02}{Treat people, living things, the environment and other people's data with respect}
\defskill{PC6}{Collect data and measure}{Blue}{VC2S8I02}{Take each measurement at the right time, with its unit, and record it straight away}
\defskill{PC7}{Use and select equipment}{Blue}{VC2S8I03}{Name each piece of equipment, use it correctly and choose the right piece for the job}
\defskill{PC8}{Use digital tools}{Blue}{VC2S8I03}{Use stopwatches, data loggers, sensors and spreadsheets to collect and record data}
% Processing, modelling and analysing — VC2S8I04, VC2S8I05
\defskill{PA1}{Tables}{GoldDark}{VC2S8I04}{Title, Headings, Units in the headings only, Decimals kept the same (THUD), with the independent variable in the first column}
\defskill{PA2}{Graphing}{GoldDark}{VC2S8I04}{The right type of graph, a title, axes labelled with units, the independent variable on the x-axis, an even scale and accurate points, then join the points or draw a line of best fit}
\defskill{PA3}{Mathematical relationships}{GoldDark}{VC2S8I04}{Process data with calculations: differences, averages, rates and formulas}
\defskill{PA4}{Creating models}{GoldDark}{VC2S8I04}{Build a physical, drawn or mathematical model that represents something real}
\defskill{PA5}{Patterns and trends}{GoldDark}{VC2S8I05}{Describe how the dependent variable changes as the independent variable changes, quoting data}
\defskill{PA6}{Correlation and causality}{GoldDark}{VC2S8I05}{Say whether two variables change together, and whether there is evidence that one causes the other}
\defskill{PA7}{Identify anomalies}{GoldDark}{VC2S8I05}{Spot a result that does not fit the pattern and suggest why it happened}
\defskill{PA8}{Interpreting models}{GoldDark}{VC2S8I05}{Use a model to explain an observation, and say what the model shows well and what it leaves out}
% Evaluating — VC2S8I06, VC2S8I07
\defskill{EV1}{Results against hypothesis}{OrangeDark}{VC2S8I06}{State the general trend and whether it supports the hypothesis (supported, never proven), with data to back the trend}
\defskill{EV2}{Compare with known concepts}{OrangeDark}{VC2S8I06}{Compare your results with accepted values and scientific ideas}
\defskill{EV3}{Error sources}{OrangeDark}{VC2S8I06}{Name what could have made results less accurate or less precise, and suggest how to improve the method}
\defskill{EV4}{Quality of secondary sources}{OrangeDark}{VC2S8I07}{Judge whether information from a website, book or video can be trusted}
\defskill{EV5}{Questions for further investigation}{OrangeDark}{VC2S8I07}{Pose a new question that follows from your results}
% Communicating — VC2S8I08
\defskill{CO1}{Scientific writing}{PurpleDark}{VC2S8I08}{Write in the third person, legibly, with correct spelling and grammar, short and to the point, using scientific terms}
\defskill{CO2}{Presentation formats}{PurpleDark}{VC2S8I08}{Choose the right format to show findings: a table, graph, scientific diagram, poster, report or talk}

% Skills focus panel at the top of each experiment:
%   \begin{skillsfocus} \focus{PC4}{what you do with it here} … \end{skillsfocus}
\newtcolorbox{skillsfocus}{
  enhanced, breakable=false, arc=2.5mm, boxrule=0.6pt,
  colback=white, colframe=NavyBlue!55,
  fonttitle=\headingfont\bfseries, coltitle=white, colbacktitle=NavyBlue,
  title={Skills focus --- the science skills this experiment builds},
  left=4mm, right=3.5mm, top=2mm, bottom=2mm}
\newcommand{\focus}[2]{\par\noindent\makebox[1.25cm][l]{\skc{#1}}\parbox[t]{\dimexpr\linewidth-1.25cm\relax}{\small\textbf{\skname{#1}}\,---\,#2}\par\vspace{2pt}}

% Learning intentions panel: \begin{lipanel} \inquiryli{…} \understandli{…} \sclist{…} \end{lipanel}
\newtcolorbox{lipanel}{
  enhanced, breakable=false, arc=2.5mm, boxrule=0.5pt,
  colback=LightBlue, colframe=Blue!40,
  fonttitle=\headingfont\bfseries, coltitle=MidBlue, title={Learning intentions},
  left=4mm, right=3.5mm, top=2mm, bottom=2mm}
\newcommand{\inquiryli}[1]{\par\noindent{\headingfont\bfseries\color{NavyBlue}Science Inquiry:}\ #1\par\vspace{2pt}}
\newcommand{\understandli}[1]{\par\noindent{\headingfont\bfseries\color{NavyBlue}Science Understanding:}\ #1\par\vspace{2pt}}
\newcommand{\sctitle}{\par\noindent{\headingfont\bfseries\color{NavyBlue}Success criteria}\par}

% Technique spotlight: \begin{technique}{title}{PC4} … \end{technique}
\newtcolorbox{technique}[2]{
  enhanced, breakable, arc=2.5mm, boxrule=0.6pt,
  colback=white, colframe=Blue!55,
  fonttitle=\headingfont\bfseries, coltitle=MidBlue, colbacktitle=LightBlue,
  title={Technique: #1\hfill\sk{#2}},
  left=4mm, right=3.5mm, top=2.5mm, bottom=2.5mm}

% Experiment divider: \experimentdiv{title}{summary}{lesson number}
% Steps the section counter, so tasks, worked examples and figures are N.M.
\newcommand{\experimentdiv}[3]{%
  \clearpage
  \refstepcounter{section}%
  \pdfbookmark[1]{Experiment \thesection{} — #1}{exp.\thesection}%
  \begin{tcolorbox}[
    enhanced, breakable=false, arc=3.5mm,
    colback=NavyBlue, colframe=NavyBlue, boxrule=0pt,
    sidebyside, sidebyside align=center,
    lefthand width=2.6cm,
    left=2mm, right=4mm, top=4mm, bottom=4mm]
    \begin{center}
      {\headingfont\bfseries\color{white!60!NavyBlue}\scriptsize EXPERIMENT}\\[-2pt]
      {\headingfont\bfseries\color{white!35!NavyBlue}\fontsize{46}{46}\selectfont \thesection}
    \end{center}
    \tcblower
    {\headingfont\bfseries\color{white}\LARGE #1}\\[3pt]
    {\itshape\color{white!75!NavyBlue}\small #2}\\[3pt]
    {\headingfont\color{Sky}\footnotesize Lesson #3}
  \end{tcolorbox}%
}

% Part bar inside an experiment. \exppart{title}{colour} starts a new page
% (new parts always start on a new page); \exppart*{…}{…} follows straight on.
\NewDocumentCommand{\exppart}{s m m}{%
  \par\IfBooleanTF{#1}{\needspace{10\baselineskip}\vspace{8pt}}{\clearpage}%
  \begin{tcolorbox}[enhanced, breakable=false, arc=2mm, boxrule=0pt,
    colback=#3!14, borderline west={4pt}{0pt}{#3},
    left=4mm, right=3mm, top=1.6mm, bottom=1.6mm]
  {\headingfont\bfseries\color{#3!60!black}\large #2}
  \end{tcolorbox}\vspace{2pt}}

% Skill check at the end of each experiment: students tick one column per row.
%   \begin{skillcheck} \scrow{PC4}{I can …} … \end{skillcheck}
\newcommand{\tickbox}{\tikz[baseline=-0.6ex]\draw[DarkGrey, line width=0.7pt, rounded corners=1pt] (0,-0.17) rectangle (0.34,0.17);}
\newenvironment{skillcheck}{%
  \par\noindent Read each statement. Tick the column that best describes you \emph{today}.\par
  \begin{ctab}\renewcommand{\arraystretch}{1.25}%
  \begin{tabular}{|W{3.9cm}|W{7.0cm}|Z{1.45cm}|Z{1.45cm}|Z{1.45cm}|}
  \hline\thead{Skill} & \thead{Success criterion} & \thead{Not yet} & \thead{Nearly} & \thead{Got it} \\ \hline}
  {\end{tabular}\end{ctab}}
\newcommand{\scrow}[2]{\rule{0pt}{1.1cm}\skc{#1}\enspace{\small\skname{#1}} & #2 & \tickbox & \tickbox & \tickbox \\ \hline}

% Write-in label under a drawing: \nameline{label}
\newcommand{\nameline}[1]{\par\noindent{\small #1}\ \leaders\hbox{\rule[-2pt]{1pt}{0.4pt}}\hfill\null\par}

% ── Equipment drawn as scientific diagrams (2D, side-on, outline only) ──────
% Each \eq... macro draws at the origin inside a tikzpicture. \eqpic[scale]{\eqBeaker}
% puts one on its own. The same shapes are used everywhere, so the booklet
% models exactly the drawings students are asked to make.
\tikzset{
  eq/.style={line width=0.9pt, draw=Ink, line cap=round, line join=round},
  eqlab/.style={draw=Ink, line width=0.4pt},
  eqtext/.style={font=\footnotesize, text=Ink, inner sep=1.5pt},
}
\newcommand{\eqpic}[2][1]{\begin{tikzpicture}[scale=#1, baseline=(current bounding box.center)]#2\end{tikzpicture}}
\newcommand{\eqBeaker}{\draw[eq] (-0.22,2.62) -- (0,2.45) -- (0,0) -- (1.8,0) -- (1.8,2.5);}
\newcommand{\eqFlask}{\draw[eq] (0.62,2.9) -- (0.65,2.8) -- (0.65,1.9) -- (0,0) -- (1.9,0) -- (1.25,1.9) -- (1.25,2.8) -- (1.28,2.9);}
\newcommand{\eqTestTube}{\draw[eq] (-0.06,3.0) -- (0,2.9) -- (0,0.3) arc[start angle=180,end angle=360,radius=0.3] -- (0.6,2.9) -- (0.66,3.0);}
\newcommand{\eqCylinder}{%
  \draw[eq] (-0.15,3.6) -- (0,3.5) -- (0,0.15) (0.6,0.15) -- (0.6,3.5);
  \draw[eq] (-0.35,0) rectangle (0.95,0.15);
  \foreach \y in {0.6,1.1,...,3.2} \draw[Ink, line width=0.5pt] (0.6,\y) -- (0.4,\y);}
\newcommand{\eqFunnel}{\draw[eq] (0,2.2) -- (0.75,1.0) -- (0.75,0) (1.85,2.2) -- (1.1,1.0) -- (1.1,0);}
\newcommand{\eqBasin}{\draw[eq] (0,0.9) .. controls (0.15,0.02) and (2.25,0.02) .. (2.4,0.9);}
\newcommand{\eqTripod}{%
  \draw[eq] (-0.15,2.0) rectangle (2.55,2.12);
  \draw[eq] (0.25,2.0) -- (0,0) (2.15,2.0) -- (2.4,0) (1.2,2.0) -- (1.2,0.15);}
\newcommand{\eqBunsen}{%
  \draw[eq] (0,0) rectangle (1.6,0.22);
  \draw[eq] (0.58,0.85) -- (0.58,2.3) (1.02,0.85) -- (1.02,2.3);
  \draw[eq] (0.5,0.22) rectangle (1.1,0.85);
  \fill[Ink] (0.72,0.42) rectangle (0.88,0.62);
  \draw[eq] (-0.65,0.06) rectangle (0,0.16);}
\newcommand{\eqThermometer}{%
  \draw[eq] (0.25,0.25) circle (0.2);
  \draw[eq] (0.15,0.43) -- (0.15,3.9) arc[start angle=180,end angle=0,radius=0.1] -- (0.35,0.43);
  \draw[Ink, line width=1.3pt] (0.25,0.45) -- (0.25,1.9);}
\newcommand{\eqStand}{%
  \draw[eq] (0,0) rectangle (2.3,0.2);
  \draw[eq] (0.35,0.2) -- (0.35,4.0) (0.5,0.2) -- (0.5,4.0);
  \draw[eq] (0.25,2.75) rectangle (0.75,3.2);
  \draw[eq] (0.75,2.92) -- (1.6,2.92) (0.75,3.03) -- (1.6,3.03);
  \draw[eq] (1.6,2.7) -- (1.6,3.25) (1.6,2.7) -- (1.95,2.7) (1.6,3.25) -- (1.95,3.25);}
\newcommand{\eqWatchGlass}{\draw[eq] (0,0.35) .. controls (0.5,-0.05) and (1.7,-0.05) .. (2.2,0.35);}
\newcommand{\eqTongs}{%
  \draw[eq] (0,0.1) -- (2.0,0.45) .. controls (2.3,0.5) and (2.45,0.35) .. (2.6,0.15)
            (0,0.0) -- (2.0,-0.35) .. controls (2.3,-0.4) and (2.45,-0.25) .. (2.6,-0.05);}

% A label line from a point on a drawing to text at the side: \eqlab{(x,y)}{(tx,ty)}{text}
\newcommand{\eqlab}[3]{\draw[eqlab] #1 -- #2 node[eqtext, anchor=west] {#3};}
\newcommand{\eqlabl}[3]{\draw[eqlab] #1 -- #2 node[eqtext, anchor=east] {#3};}

% ── Reading scales ──────────────────────────────────────────────────────────
% \readcyl{lowest}{highest}{labelled step}{small divisions per label}{reading}
% A close-up of a measuring cylinder's scale, the water drawn to the reading:
% the bottom of the meniscus sits exactly on the reading.
\pgfkeys{/pgf/number format/.cd, fixed, precision=1}
\newcommand{\readcyl}[5]{%
\begin{tikzpicture}[baseline=(current bounding box.south)]
  \pgfmathsetmacro\rcH{5.0}
  \pgfmathsetmacro\rcY{(#5-#1)*\rcH/(#2-#1)}
  \pgfmathtruncatemacro\rcN{round((#2-#1)/#3*#4)}
  \fill[black!10] (0,-0.3) -- (0,\rcY+0.16) .. controls (0.45,\rcY-0.05) and (1.15,\rcY-0.05) .. (1.6,\rcY+0.16) -- (1.6,-0.3) -- cycle;
  \draw[Ink, line width=0.9pt] (0,\rcY+0.16) .. controls (0.45,\rcY-0.05) and (1.15,\rcY-0.05) .. (1.6,\rcY+0.16);
  \foreach \k in {0,...,\rcN}{
    \pgfmathsetmacro\yy{\k*\rcH/\rcN}
    \pgfmathtruncatemacro\ismaj{mod(\k,#4)==0 ? 1 : 0}
    \ifnum\ismaj=1
      \draw[Ink, line width=0.75pt] (0,\yy) -- (0.6,\yy);
      \pgfmathsetmacro\lv{#1+\k/#4*#3}
      \node[anchor=west, font=\scriptsize\bfseries, inner sep=1pt, fill=white, fill opacity=0.8, text opacity=1] at (0.66,\yy) {\pgfmathprintnumber{\lv}};
    \else
      \draw[Ink, line width=0.4pt] (0,\yy) -- (0.32,\yy);
    \fi}
  \draw[eq] (0,-0.3) -- (0,\rcH+0.35) (1.6,-0.3) -- (1.6,\rcH+0.35);
  \node[font=\scriptsize, anchor=north] at (0.8,-0.35) {scale in mL};
\end{tikzpicture}}
% \readtherm{lowest}{highest}{labelled step}{small divisions per label}{reading}
\newcommand{\readtherm}[5]{%
\begin{tikzpicture}[baseline=(current bounding box.south)]
  \pgfmathsetmacro\rtH{5.0}
  \pgfmathsetmacro\rtY{(#5-#1)*\rtH/(#2-#1)}
  \pgfmathtruncatemacro\rtN{round((#2-#1)/#3*#4)}
  \draw[eq, fill=white] (0,-0.45) circle (0.3);
  \fill[black!75] (0,-0.45) circle (0.22);
  \fill[black!75] (-0.07,-0.3) rectangle (0.07,\rtY);
  \draw[eq] (-0.16,-0.2) -- (-0.16,\rtH+0.3) arc[start angle=180,end angle=0,radius=0.16] -- (0.16,-0.2);
  \foreach \k in {0,...,\rtN}{
    \pgfmathsetmacro\yy{\k*\rtH/\rtN}
    \pgfmathtruncatemacro\ismaj{mod(\k,#4)==0 ? 1 : 0}
    \ifnum\ismaj=1
      \draw[Ink, line width=0.75pt] (0.25,\yy) -- (0.75,\yy);
      \pgfmathsetmacro\lv{#1+\k/#4*#3}
      \node[anchor=west, font=\scriptsize\bfseries, inner sep=1pt] at (0.8,\yy) {\pgfmathprintnumber{\lv}};
    \else
      \draw[Ink, line width=0.4pt] (0.25,\yy) -- (0.52,\yy);
    \fi}
  \node[font=\scriptsize, anchor=north] at (0.3,-0.8) {scale in \textdegree C};
\end{tikzpicture}}
% \balancefig{reading}: an electronic balance showing a reading
\newcommand{\balancefig}[1]{%
\begin{tikzpicture}[baseline=(current bounding box.south)]
  \draw[eq, rounded corners=3pt] (0,0) rectangle (4.2,1.35);
  \draw[eq] (1.9,1.35) -- (1.9,1.55) (2.3,1.35) -- (2.3,1.55);
  \draw[eq, rounded corners=1pt] (0.5,1.55) rectangle (3.7,1.72);
  \draw[eq, fill=LightGrey] (0.35,0.35) rectangle (2.75,1.02);
  \node[font=\headingfont\bfseries\large, anchor=east] at (2.65,0.68) {#1};
  \draw[eq] (3.3,0.68) circle (0.2);
  \node[font=\tiny\bfseries] at (3.3,0.3) {TARE};
\end{tikzpicture}}
% Bunsen burner with a flame: \bunsenflame{safety|heating}
\newcommand{\bunsenflame}[1]{%
\begin{tikzpicture}[baseline=(current bounding box.south)]
  \def\fl{#1}\def\safe{safety}%
  \ifx\fl\safe
    \fill[black!14] (0.8,2.3) .. controls (0.05,3.0) and (0.45,4.3) .. (0.85,5.0) .. controls (1.05,4.4) and (1.55,3.1) .. (0.8,2.3);
    \draw[Ink, line width=0.6pt] (0.8,2.3) .. controls (0.05,3.0) and (0.45,4.3) .. (0.85,5.0) .. controls (1.05,4.4) and (1.55,3.1) .. (0.8,2.3);
  \else
    \fill[white] (0.8,2.3) .. controls (0.45,2.8) and (0.6,3.5) .. (0.8,3.9) .. controls (1.0,3.5) and (1.15,2.8) .. (0.8,2.3);
    \draw[Ink, line width=0.6pt] (0.8,2.3) .. controls (0.45,2.8) and (0.6,3.5) .. (0.8,3.9) .. controls (1.0,3.5) and (1.15,2.8) .. (0.8,2.3);
    \fill[black!30] (0.8,2.3) .. controls (0.65,2.55) and (0.7,2.85) .. (0.8,3.0) .. controls (0.9,2.85) and (0.95,2.55) .. (0.8,2.3);
  \fi
  \draw[eq] (0,0) rectangle (1.6,0.22);
  \draw[eq] (0.58,0.85) -- (0.58,2.3) (1.02,0.85) -- (1.02,2.3);
  \draw[eq] (0.5,0.22) rectangle (1.1,0.85);
  \ifx\fl\safe\draw[Ink, line width=0.5pt, densely dotted] (0.72,0.42) rectangle (0.88,0.62);\else\fill[Ink] (0.72,0.42) rectangle (0.88,0.62);\fi
  \draw[eq] (-0.65,0.06) rectangle (0,0.16);
\end{tikzpicture}}
% ══════════════════════════════════════════════════════════════════════════════
%  WATER CYCLE AND RESOURCES — COVER
%  The unit's cover art (covers/water-resources.jpg), full bleed. The A5 lesson
%  booklet sets \wrcoverfile one folder deeper.
% ══════════════════════════════════════════════════════════════════════════════
\providecommand{\wrcoverfile}{../../covers/water-resources.jpg}
\newcommand{\wrcover}{%
\begin{tikzpicture}[remember picture, overlay]
  \node[anchor=north] at (current page.north) {\includegraphics[width=\paperwidth]{\wrcoverfile}};
\end{tikzpicture}}
% ══════════════════════════════════════════════════════════════════════════════
%  2027 WORKBOOK STYLE (Year 7 workbook redesign, Nov 2026) — A4 version
%  The look of the A5 trial (space/booklet/a5-test) on an A4 page:
%    · colour cover from the unit's cover art, black-and-white inside
%    · headings in Outfit (the site's heading face), body in Lato
%    · each experiment opens with a band: number, title, ink symbol, heavy rule
%    · boxes told apart by frame and icon, never by colour:
%        I do ●  solid · We do ●● dashed · You do ●●● heavy
%        Key idea / technique ◆ thin · Above level ▲ double · Country ◎ dotted
%    · skill tags are ink tabs carrying the code (QP1, PC4 …) plus the name
%  Everything above this block still names the old colours; they are mapped to
%  ink here, so the body never needs to know which look it is printed in.
% ══════════════════════════════════════════════════════════════════════════════

\usetikzlibrary{patterns}
% ── Fonts ───────────────────────────────────────────────────────────────────
% Bold is Outfit SemiBold, the weight the site uses; Outfit has no italic,
% so a slanted one stands in.
\renewfontfamily\headingfont{Outfit}[
  Path=./fonts/, Extension=.ttf, Ligatures=TeX,
  UprightFont=Outfit-Regular, BoldFont=Outfit-SemiBold,
  ItalicFont=Outfit-Regular, ItalicFeatures={FakeSlant=0.18},
  BoldItalicFont=Outfit-SemiBold, BoldItalicFeatures={FakeSlant=0.18} ]

% ── Black and white inside ──────────────────────────────────────────────────
\definecolor{Ink}{HTML}{111111}
\colorlet{NavyBlue}{Ink}\colorlet{MidBlue}{Ink}\colorlet{TealGreen}{Ink}
\colorlet{GoldDark}{Ink}\colorlet{PurpleDark}{Ink}\colorlet{OrangeDark}{Ink}
\colorlet{TableHead}{Ink}
\colorlet{Blue}{Ink}\colorlet{Teal}{Ink}\colorlet{Amber}{Ink}\colorlet{Coral}{Ink}
\colorlet{Violet}{Ink}\colorlet{Sky}{Ink}
\colorlet{LightBlue}{white}\colorlet{LightGreen}{white}\colorlet{GoldLight}{white}
\colorlet{PurpleLight}{white}\colorlet{OrangeLight}{white}\colorlet{RowAlt}{white}
\colorlet{LightGrey}{white}\colorlet{Paper}{white}
\colorlet{DarkGrey}{black!72}\colorlet{MidGrey}{black!45}\colorlet{Line}{black!45}
\arrayrulecolor{black!60}

% ── Page style ──────────────────────────────────────────────────────────────
\fancypagestyle{fancy}{%
  \fancyhf{}%
  \renewcommand{\headrulewidth}{0.4pt}%
  \renewcommand{\headrule}{\hbox to\headwidth{\color{Ink}\leaders\hrule height \headrulewidth\hfill}}%
  \fancyhead[L]{\headingfont\footnotesize \unitname\ \textbullet\ Experiment Booklet}%
  \fancyhead[R]{\headingfont\footnotesize\leftmark}%
  \fancyfoot[C]{\headingfont\footnotesize\thepage}}
\fancypagestyle{plain}{\fancyhf{}\renewcommand{\headrulewidth}{0pt}%
  \fancyfoot[C]{\headingfont\footnotesize\thepage}}
\fancypagestyle{opener}{\fancyhf{}\renewcommand{\headrulewidth}{0pt}%
  \fancyfoot[C]{\headingfont\footnotesize\thepage}}
\pagestyle{fancy}
\markboth{}{}

\titleformat{\subsubsection}{\headingfont\normalsize\bfseries}{}{0em}{}

% ── Box system: told apart by frame and icon, never by colour ───────────────
\newcommand{\icodots}[1]{\tikz[baseline=-0.6ex]{\foreach \k in {1,...,#1}{\fill[white] (\k*1.9mm,0) circle (0.7mm);}}}
\newcommand{\icokey}{\tikz[baseline=-0.6ex]{\fill[white] (0,0.9mm)--(0.9mm,0)--(0,-0.9mm)--(-0.9mm,0)--cycle;}}
\newcommand{\icoup}{\tikz[baseline=-0.6ex]{\fill[white] (-1mm,-0.8mm)--(1mm,-0.8mm)--(0,1mm)--cycle;}}
\newcommand{\icocountry}{\tikz[baseline=-0.6ex]{\draw[white,line width=0.5pt] (0,0) circle (1mm); \fill[white] (0,0) circle (0.45mm);}}
\newcommand{\icoskill}{\tikz[baseline=-0.6ex]{\fill[white] (-0.85mm,-0.85mm) rectangle (0.85mm,0.85mm);}}
\newcommand{\tagicon}[1]{\tikz[baseline=-0.6ex]\node[fill=Ink, rounded corners=0.8mm, inner xsep=1.2mm, inner ysep=1mm]{#1};}
\tcbset{wbbase/.style={enhanced, colback=white, arc=1.2mm,
    left=3.5mm, right=3.5mm, top=4.2mm, bottom=2.6mm, before skip=11pt, after skip=8pt,
    attach boxed title to top left={xshift=3mm, yshift=-2.4mm},
    boxed title style={colback=Ink, colframe=Ink, arc=0.8mm, boxrule=0pt,
                       left=1.6mm, right=1.8mm, top=0.6mm, bottom=0.6mm},
    fonttitle=\headingfont\bfseries\footnotesize, coltitle=white}}

% I do: the worked example (solid frame, one dot)
\renewtcolorbox{workedexbox}[1]{wbbase, breakable, colframe=Ink, boxrule=1pt,
  title={\icodots{1}\ I do \textbullet\ #1}}
\renewenvironment{workedex}[1]
  {\refstepcounter{wenum}\label{we:\thewenum}\begin{workedexbox}{Worked example \thewenum\ --- #1}}
  {\end{workedexbox}}
% We do: straight after its worked example (dashed frame, two dots)
\renewtcolorbox{wedobox}[1]{wbbase, breakable=false, frame hidden,
  borderline={1.1pt}{0pt}{Ink, dash pattern=on 4pt off 2.6pt},
  title={\icodots{2}\ #1}}
\renewenvironment{wedo}[1]
  {\par\begin{wedobox}{We do \thewenum\ --- #1}\label{wd:\thewenum}\ignorespaces}
  {\end{wedobox}}
% You do: a set of practice tasks on your own (heavy frame, three dots)
\newtcolorbox{youdobox}[1][Check your understanding]{wbbase, breakable=false,
  colframe=Ink, boxrule=2.4pt, top=5mm,
  title={\icodots{3}\ You do \textbullet\ #1}}
% Key idea, technique, safety, word banks (thin frame, diamond)
\renewtcolorbox{infobox}[1]{wbbase, breakable, colframe=Ink, boxrule=0.5pt,
  title={\icokey\ #1}}
\renewtcolorbox{scaffold}[1]{wbbase, breakable, colframe=Ink, boxrule=0.5pt,
  title={\icokey\ #1}}
\renewtcolorbox{technique}[2]{wbbase, breakable, colframe=Ink, boxrule=0.5pt,
  title={\icokey\ Technique \textbullet\ #1\enspace\textbar\enspace #2 \skname{#2}}}
% Above level (double frame, triangle)
\renewtcolorbox{yr8box}{wbbase, breakable, frame hidden,
  borderline={0.6pt}{0pt}{Ink}, borderline={0.6pt}{1.6pt}{Ink},
  title={\icoup\ Going further \textbullet\ Above level}}
% Country (dotted frame, ring)
\newtcolorbox{countrybox}[1]{wbbase, breakable, frame hidden,
  borderline={1.4pt}{0pt}{Ink, line cap=round, dash pattern=on 0pt off 2.4pt},
  title={\icocountry\ #1}}
% Learning intentions and skills focus (thin frame; square for skills)
\renewtcolorbox{lipanel}{wbbase, breakable=false, colframe=Ink, boxrule=0.5pt,
  title={\icokey\ Learning intentions}}
\renewtcolorbox{skillsfocus}{wbbase, breakable=false, colframe=Ink, boxrule=1pt,
  title={\icoskill\ Skills focus \textbullet\ the science skills this experiment builds},
  before upper={\begin{multicols}{2}\raggedcolumns}, after upper={\end{multicols}}}

% ── Skill tags: an ink tab with the code, then the name ─────────────────────
\makeatletter
\renewcommand{\skc}[1]{%
  \tikz[baseline=(c.base)]\node[rounded corners=0.8mm, fill=Ink, text=white, inner xsep=1.5mm, inner ysep=0.7mm,
    font=\headingfont\bfseries\scriptsize] (c) {#1};}
\renewcommand{\sk}[1]{\mbox{\skc{#1}\,{\headingfont\bfseries\footnotesize\skname{#1}}}}
\makeatother
\renewcommand{\focus}[2]{\par\noindent\makebox[1.1cm][l]{\skc{#1}}\parbox[t]{\dimexpr\linewidth-1.1cm\relax}{\footnotesize\RaggedRight\textbf{\skname{#1}}\,---\,#2}\par\vspace{3pt}}

% ── Experiment opener band ──────────────────────────────────────────────────
% \experimentdiv{title}{summary}{lesson}: number, title and the experiment's ink
% symbol (\expsymbol{N}), then a heavy rule. Steps the section counter, so
% tasks, worked examples and figures are N.M.
\newcommand{\expword}{Experiment}\newcommand{\expcaps}{EXPERIMENT}\newcommand{\symid}{\thesection}
\renewcommand{\experimentdiv}[3]{%
  \clearpage\thispagestyle{opener}%
  \refstepcounter{section}\markboth{\expword\ \thesection\ \textbullet\ #1}{}%
  \pdfbookmark[1]{\expword\ \thesection{} — #1}{exp.\thesection}%
  \noindent\begin{minipage}[b]{\dimexpr\linewidth-46mm\relax}
    {\headingfont\bfseries\footnotesize\addfontfeatures{LetterSpace=14}\unitcaps\ \textbullet\ \expcaps}\par\vspace{1mm}
    {\headingfont\bfseries\fontsize{40}{40}\selectfont\thesection}\par\vspace{1mm}
    {\headingfont\bfseries\fontsize{19}{22}\selectfont\RaggedRight #1\par}\vspace{1mm}
    {\small\itshape #2}\par\vspace{1mm}
    {\headingfont\footnotesize Lesson #3}
  \end{minipage}\hfill
  \begin{minipage}[b]{44mm}\centering\resizebox{!}{27mm}{\expsymbol{\symid}}\end{minipage}\par
  \vspace{2mm}\noindent\rule{\linewidth}{1.6pt}\par\vspace{1mm}}
\renewcommand{\lessoncodes}[1]{{\small\itshape VC 2.0: #1}\par\vspace{2pt}}

% Part heading inside an experiment: every part starts a new page; the
% starred form follows straight on.
\RenewDocumentCommand{\exppart}{s m m}{%
  \par\IfBooleanTF{#1}{\needspace{10\baselineskip}\vspace{8pt}}{\clearpage}%
  {\headingfont\bfseries\LARGE #2}\par\vspace{1mm}\noindent\rule{\linewidth}{0.6pt}\par\vspace{2mm}}

% Key vocabulary: a heading like the trial's, then the write-in table
\renewcommand{\vocabhead}{%
  \par\vspace{4pt}{\headingfont\bfseries\Large Key vocabulary}\par\vspace{1pt}%
  \noindent{\small Copy the meaning of each word from the slide. The first one has been done for you.}\par}

% Table headings: bold Outfit on white, the grid in ink
\renewcommand{\thead}[1]{\headingfont\bfseries #1}
\renewcommand{\tickbox}{\tikz[baseline=-0.6ex]\draw[Ink, line width=0.7pt, rounded corners=1pt] (0,-0.17) rectangle (0.34,0.17);}

% Glossary rows: \glossrow{word}{label}
\newcommand{\glossrow}[2]{\rule{0pt}{5mm}\textbf{#1} & \pageref{#2} \\ \hline}

% ── Ink symbols for the experiment openers ───────────────────────────────────
% One diagrammatic idea per experiment, drawn in ink from the same 2D
% equipment shapes the booklet teaches.
\tikzset{sym/.style={line width=1.5pt, draw=Ink, line cap=round, line join=round}}
\newcommand{\expsymbol}[1]{%
  \begin{tikzpicture}[x=1cm,y=1cm]
  \ifcase#1
  \or
    % 1 Water cycle in a bag: a sealed bag with water at the bottom, drops on the inside, the sun beside it
    \draw[sym] (0.2,0) rectangle (2.6,3.2);
    \draw[Ink, line width=1pt] (0.2,2.9) -- (2.6,2.9);
    \fill[black!14] (0.25,0.05) rectangle (2.55,0.75);
    \draw[Ink, line width=1pt] (0.25,0.75) .. controls (0.8,0.85) and (1.9,0.65) .. (2.55,0.75);
    \foreach \p in {(0.6,2.3),(1.1,2.55),(1.7,2.35),(2.2,2.6),(0.9,1.9),(2.0,1.95)}{\fill[Ink] \p circle (0.07);}
    \draw[Ink, line width=1.2pt, -{Stealth[length=5pt]}] (1.4,1.0) -- (1.4,1.7);
    \draw[sym, fill=white] (3.6,2.4) circle (0.45);
    \foreach \a in {0,45,...,315}{\draw[Ink, line width=1.4pt] ($(3.6,2.4)+(\a:0.6)$) -- ($(3.6,2.4)+(\a:0.85)$);}
  \or
    % 2 Household energy audit: an appliance label with a star rating, and a power plug
    \draw[sym] (0,0) rectangle (2.2,3.1);
    \foreach \x in {0.3,0.7,1.1,1.5,1.9}{\draw[Ink, line width=1pt, fill=white] ($(\x,2.5)+(90:0.18)$) -- ($(\x,2.5)+(162:0.18)$) -- ($(\x,2.5)+(234:0.18)$) -- ($(\x,2.5)+(306:0.18)$) -- ($(\x,2.5)+(18:0.18)$) -- cycle;}
    \fill[Ink] ($(0.3,2.5)+(90:0.18)$) -- ($(0.3,2.5)+(162:0.18)$) -- ($(0.3,2.5)+(234:0.18)$) -- ($(0.3,2.5)+(306:0.18)$) -- ($(0.3,2.5)+(18:0.18)$) -- cycle;
    \fill[Ink] ($(0.7,2.5)+(90:0.18)$) -- ($(0.7,2.5)+(162:0.18)$) -- ($(0.7,2.5)+(234:0.18)$) -- ($(0.7,2.5)+(306:0.18)$) -- ($(0.7,2.5)+(18:0.18)$) -- cycle;
    \fill[Ink] ($(1.1,2.5)+(90:0.18)$) -- ($(1.1,2.5)+(162:0.18)$) -- ($(1.1,2.5)+(234:0.18)$) -- ($(1.1,2.5)+(306:0.18)$) -- ($(1.1,2.5)+(18:0.18)$) -- cycle;
    \draw[Ink, line width=1pt] (0.3,1.85) -- (1.9,1.85) (0.3,1.45) -- (1.5,1.45);
    \node[font=\headingfont\bfseries\large] at (1.1,0.75) {kWh};
    \draw[sym] (3.0,1.3) rectangle (3.9,2.2);
    \draw[sym] (3.2,2.2) -- (3.2,2.7) (3.7,2.2) -- (3.7,2.7);
    \draw[sym] (3.45,1.3) .. controls (3.45,0.6) and (2.6,0.6) .. (2.5,0.0);
  \fi
  \end{tikzpicture}}
\begin{document}
\setlength{\parindent}{0pt}
\setlength{\parskip}{4pt}

% ---- Cover: the unit's cover art, full bleed, in colour ------------------
% The art is 2:3 and A4 is taller-than-wide by less, so it fills the width and
% loses a little off the bottom; the baked-in title at the top stays whole.
\pdfbookmark[1]{Cover}{cover}
\thispagestyle{empty}
\wrcover
\begin{tikzpicture}[remember picture, overlay]
  
  % Name box, top right, clear of the title.
  \node[anchor=north east, fill=white, draw=Ink, line width=0.6pt, rounded corners=1.5mm,
        inner xsep=2.6mm, inner ysep=2.2mm, text width=40mm, font=\headingfont\small]
    at ([xshift=-6mm, yshift=-6mm]current page.north east)
    {\textbf{Name}\par\vspace{5mm}\hrule\vspace{2.6mm}\textbf{Class}\par\vspace{5mm}\hrule};
  % What this book is: the cover art says Workbook, so name the booklet here.
  \node[anchor=south east, fill=white, draw=Ink, line width=0.6pt, rounded corners=1.5mm,
        inner xsep=3.2mm, inner ysep=2.4mm, font=\headingfont\bfseries\large, align=right]
    at ([xshift=-6mm, yshift=8mm]current page.south east)
    {Experiment Booklet\\[-1pt]{\normalfont\headingfont\small Experiments 1--2}};
\end{tikzpicture}
\mbox{}
\clearpage

% ---- Inside cover ---------------------------------------------------------
\thispagestyle{plain}
{\headingfont\bfseries\fontsize{26}{30}\selectfont Water Cycle and Resources}\par\vspace{1mm}
{\headingfont\Large Experiment Booklet \textbullet\ Year 7}\par\vspace{1mm}
{\small\itshape VC 2.0: VC2S8U09 \textbullet\ VC2S8U16 \textbullet\ VC2S8I01--I08 \textbullet\ VC2S8H04}\par\vspace{6mm}
\renewcommand{\arraystretch}{2.2}
\begin{tabular}{@{}l p{130mm}@{}}
\textbf{Name} & \hrulefill \\ \textbf{Class} & \hrulefill \\ \textbf{Teacher} & \hrulefill \\ \textbf{Lab group} & \hrulefill \\
\end{tabular}
\renewcommand{\arraystretch}{1}

\begin{countrybox}{Acknowledgement of Country}
This booklet was created on the unceded Country of the Wurundjeri Woi-wurrung people of the Kulin Nation. We acknowledge Elders past, present and emerging. Aboriginal and Torres Strait Islander peoples have observed, tested and passed on knowledge of Country for tens of thousands of years: science in the way this booklet teaches it, by careful observation, measurement and sharing what was found.
\end{countrybox}

\begin{infobox}{How this booklet works}
This booklet holds every experiment in the unit. Bring it to every practical lesson. Each experiment opens with its key vocabulary, learning intentions and skills focus, then runs \emph{Technique} $\rightarrow$ \emph{Plan and safety} $\rightarrow$ \emph{Method} $\rightarrow$ \emph{Results} $\rightarrow$ \emph{Analysis and evaluation} $\rightarrow$ \emph{Skill check}. Each kind of box has its own frame and mark:\par\smallskip
\renewcommand{\arraystretch}{1.35}
\begin{tabular}{@{}p{32mm}l@{}}
\textbf{Technique} \tagicon{\icokey} & How to do it properly (thin frame) \\
\textbf{I do} \tagicon{\icodots{1}} & Your teacher works an example through --- you follow (solid frame) \\
\textbf{We do} \tagicon{\icodots{2}} & You try one together, with help (dashed frame) \\
\textbf{You do} \tagicon{\icodots{3}} & On your own: check your understanding (heavy frame) \\
\textbf{Going further} \tagicon{\icoup} & Above level: Year 8 extension (double frame) \\
\textbf{Skills focus} \tagicon{\icoskill} & The science skills the experiment builds \\
\end{tabular}

\smallskip
Every skill comes from the faculty \textbf{Science Skills} list (the toolkit on page \pageref{sec:toolkit}). A skill always carries the same code, for example \sk{PC4} or \sk{PA1}, so you can follow it from this unit into every other unit this year.
\end{infobox}

\clearpage
{\headingfont\bfseries\LARGE Contents}\par\vspace{1mm}\noindent\rule{\linewidth}{0.6pt}\par\vspace{2mm}
\begin{ctab}
\renewcommand{\arraystretch}{1.5}
\begin{tabular}{|L{2.8cm}|L{10.2cm}|C{1.6cm}|}
  \hline\thead{Section} & \thead{Title} & \thead{Page} \\ \hline
  Skills & \hyperlink{toolkit}{The Science Skills toolkit} & \hyperlink{toolkit}{\pageref{sec:toolkit}} \\ \hline
  Skills & \hyperlink{tracker}{My skills tracker} & \hyperlink{tracker}{\pageref{sec:tracker}} \\ \hline
  Experiment 1 & \hyperlink{exp.1}{Water Cycle in a Bag (Lesson 2)} & \hyperlink{exp.1}{\pageref{sec:1}} \\ \hline
  Experiment 2 & \hyperlink{exp.2}{Household Energy Audit (Lesson 6)} & \hyperlink{exp.2}{\pageref{sec:2}} \\ \hline
  Glossary & \hyperlink{glossary}{Key vocabulary, with the page each word first appears on} & \hyperlink{glossary}{\pageref{sec:glossary}} \\ \hline
\end{tabular}
\end{ctab}

% ══════════════════════════════════════════════════════════════════════════════
%  THE SCIENCE SKILLS TOOLKIT — faculty continuum, Years 7–8 column
% ══════════════════════════════════════════════════════════════════════════════
\clearpage
\hypertarget{toolkit}{}\pdfbookmark[1]{The Science Skills toolkit}{toolkitbm}
\markboth{The Science Skills toolkit}{}{\headingfont\bfseries\LARGE The Science Skills toolkit}\label{sec:toolkit}\par\vspace{1mm}\noindent\rule{\linewidth}{0.6pt}\par
These are the skills scientists use in every investigation. They come from the Science Inquiry strand of the Victorian Curriculum, grouped the way our faculty teaches them. Each experiment in this booklet tells you which skills it builds.

\newcommand{\tkhead}[2]{\multicolumn{3}{|l|}{\cellcolor{black!10}\rule{0pt}{4.6mm}\headingfont\bfseries #2} \\ \hline}
\newcommand{\tkrow}[1]{\skc{#1}\enspace{\small\bfseries\skname{#1}} & \skdesc{#1} & {\small\skcode{#1}} \\ \hline}
\renewcommand{\arraystretch}{1.35}
\begin{ctab}
\begin{longtable}{|W{5.0cm}|W{9.0cm}|Z{1.7cm}|}
  \hline\thead{Skill} & \thead{What it means} & \thead{Code} \\ \hline
  \endhead
  \tkhead{Teal}{Questioning and predicting}
  \tkrow{QP1}
  \tkrow{QP2}
  \tkhead{Blue}{Planning and conducting}
  \tkrow{PC1}
  \tkrow{PC2}
  \tkrow{PC3}
  \tkrow{PC4}
  \tkrow{PC5}
  \tkrow{PC6}
  \tkrow{PC7}
  \tkrow{PC8}
  \tkhead{GoldDark}{Processing, modelling and analysing}
  \tkrow{PA1}
  \tkrow{PA2}
  \tkrow{PA3}
  \tkrow{PA4}
  \tkrow{PA5}
  \tkrow{PA6}
  \tkrow{PA7}
  \tkrow{PA8}
  \tkhead{OrangeDark}{Evaluating}
  \tkrow{EV1}
  \tkrow{EV2}
  \tkrow{EV3}
  \tkrow{EV4}
  \tkrow{EV5}
  \tkhead{PurpleDark}{Communicating}
  \tkrow{CO1}
  \tkrow{CO2}
\end{longtable}
\end{ctab}
\renewcommand{\arraystretch}{1}

\begin{infobox}{Writing like a scientist \textbar\ CO1 \skname{CO1}}
Scientists write in the \textbf{third person}: say what was done, not who did it. Write \emph{``The water was heated for 10 minutes''}, not \emph{``I heated the water for 10 minutes''}. Keep sentences short, spell scientific words correctly, and always give a unit with a measurement.
\end{infobox}

% ══════════════════════════════════════════════════════════════════════════════
%  SKILLS TRACKER
% ══════════════════════════════════════════════════════════════════════════════
\clearpage
\hypertarget{tracker}{}\pdfbookmark[1]{My skills tracker}{trackerbm}
\markboth{My skills tracker}{}{\headingfont\bfseries\LARGE My skills tracker}\label{sec:tracker}\par\vspace{1mm}\noindent\rule{\linewidth}{0.6pt}\par
A dot \textbullet{} shows the experiments that build each skill. After each experiment, write how confident you are with each of its skills in its column: \textbf{1}~not yet, \textbf{2}~nearly, \textbf{3}~got it.

\newcommand{\dotc}{{\Large\textbullet}}
\newcommand{\trow}[3]{\rule[-0.36cm]{0pt}{0.92cm}\skc{#1}\enspace{\small\skname{#1}} & #2 & #3 & \\ \hline}
\begin{ctab}
\begin{tabular}{|W{6.4cm}|Z{2.2cm}|Z{2.2cm}|Z{2.6cm}|}
  \hline\thead{Skill} & \thead{Exp 1} & \thead{Exp 2} & \thead{End of unit} \\ \hline
  \trow{QP1}{\dotc}{}
  \trow{PC5}{}{\dotc}
  \trow{PC6}{\dotc}{\dotc}
  \trow{PA1}{}{\dotc}
  \trow{PA2}{}{\dotc}
  \trow{PA3}{}{\dotc}
  \trow{PA4}{\dotc}{}
  \trow{PA5}{}{\dotc}
  \trow{PA6}{}{\dotc}
  \trow{PA8}{\dotc}{}
  \trow{EV2}{\dotc}{}
  \trow{EV3}{\dotc}{}
  \trow{EV4}{}{\dotc}
  \trow{EV5}{}{\dotc}
  \trow{CO1}{}{\dotc}
  \trow{CO2}{\dotc}{}
\end{tabular}
\end{ctab}
{\small\color{DarkGrey} The other skills are built in the experiment booklets for other units.}

% ══════════════════════════════════════════════════════════════════════════════
%  EXPERIMENT 1 — WATER CYCLE IN A BAG                                  (Lesson 2)
% ══════════════════════════════════════════════════════════════════════════════
\experimentdiv{Water Cycle in a Bag}{Building a model \textbullet{} Observing over several days \textbullet{} Strengths and limitations of a model}{2}\label{sec:1}
\lessoncodes{VC2S8U09 \textbullet{} VC2S8I01, I03, I05, I06, I08}

\vocabhead
\begin{ctab}
\begin{tabular}{|SW{4.2cm}|W{11.6cm}|}
  \hline\hd & \thead{Word} & \thead{Definition (copy from the slide)} \\ \hline
  & \scaf{\textbf{Model}\label{voc:model}} & \scaf{A simpler version of something real, used to explain it or make predictions.} \\ \hline
  & \textbf{Evaporation}\label{voc:evaporation} & \\ \hline
  & \textbf{Condensation}\label{voc:condensation} & \\ \hline
  & \textbf{Precipitation}\label{voc:precipitation} & \\ \hline
  & \textbf{Limitation}\label{voc:limitation} & \\ \hline
\end{tabular}
\end{ctab}

\begin{lipanel}
\inquiryli{Make and record observations of a model over time, and evaluate the model.}
\understandli{Identify evaporation, condensation and precipitation in the model.}
\sctitle
\begin{itemize}[leftmargin=1.4em, itemsep=0pt, topsep=1pt]
  \item I can record clear observations of my model over several days.
  \item I can match each part of my model to a stage of the water cycle.
  \item I can give one strength and one limitation of the model.
\end{itemize}
\end{lipanel}

\begin{skillsfocus}
\focus{QP1}{Predict what will happen in the bag.}
\focus{PC6}{Observe the bag at each lesson.}
\focus{PA4}{Build a model of the water cycle.}
\focus{PA8}{Match the model to the real cycle.}
\focus{EV2}{Compare the sunny bag and the shaded bag.}
\focus{EV3}{Name the model's limitations.}
\focus{CO2}{Draw the set-up as a scientific diagram.}
\end{skillsfocus}

% ─── TECHNIQUE ────────────────────────────────────────────────────────────────
\exppart{Technique spotlight}{Blue}

\begin{technique}{Observing over time}{PC6}
Some experiments change slowly, so you observe them again and again. To make the observations comparable:
\begin{enumerate}[leftmargin=1.6em, itemsep=1pt, topsep=3pt]
  \item Observe at about the same time of day, from the same position.
  \item Record the date, time and weather each time.
  \item Write what you see in exact words: \emph{where}, \emph{how much}, \emph{what size}.
  \item Add a quick labelled sketch: a picture often shows what words miss.
\end{enumerate}
\end{technique}

\begin{workedex}{A precise observation}
\textbf{Vague:} \emph{The bag looked wet.}

\smallskip
\textbf{Precise:} \emph{Small clear droplets covered the top third of the bag. A few larger drops had run down the side facing the window.}
\end{workedex}

\begin{wedo}{A precise observation}
Rewrite this observation so that it is precise: \emph{There was some water at the bottom.}
\answerbox{3.5cm}
\end{wedo}

\begin{technique}{Strengths and limitations of a model}{PA8}
A \textbf{model} is useful because it shows something real in a simpler way. Every model also leaves things out or gets things wrong: these are its \textbf{limitations}. Judge a model by saying what it shows well (a strength) and what it does not show (a limitation).
\end{technique}

\begin{workedex}{A globe}
\textbf{Strength:} a globe shows that Earth is round, and where the continents and oceans are.

\smallskip
\textbf{Limitation:} it does not show weather or clouds, and its mountains are not to scale.
\end{workedex}

\begin{wedo}{A train-line map}
A map of Melbourne's train lines is a model. Give one strength and one limitation.
\answerbox{4.5cm}
\end{wedo}

\begin{youdobox}[Matching the model]
\task Each part of the bag model stands for something in the real water cycle. Complete the table. \skr{PA8}
\begin{ctab}
\begin{tabular}{|SW{5.0cm}|W{4.6cm}|W{5.4cm}|}
  \hline\hd & \thead{In the bag} & \thead{In the real world} & \thead{Stage of the water cycle} \\ \hline
  & \scaf{Water at the bottom} & \scaf{Oceans, lakes and rivers} & \scaf{Collection} \\ \hline
  & Sunlight through the window & & \\ \hline
  & Water vapour in the air of the bag & & \\ \hline
  & Droplets on the top of the bag & & \\ \hline
  & Drops running down the sides & & \\ \hline
\end{tabular}
\end{ctab}
\end{youdobox}

% ─── PLAN AND SAFETY / METHOD ────────────────────────────────────────────────
\exppart{Plan, safety and method}{Teal}

\textbf{Aim:} To model the water cycle in a sealed bag and observe evaporation, condensation and precipitation.

\needspace{10\baselineskip}
\task Predict what you will see inside a bag taped to a sunny window after one day, and inside a bag kept in a dark cupboard. Explain your prediction. \skr{QP1}
\answerbox{5cm}

\subsubsection*{Equipment}
Resealable plastic bag (about 20\,cm $\times$ 25\,cm) \textbullet{} 100\,mL measuring cylinder \textbullet{} water and blue food colouring \textbullet{} permanent marker \textbullet{} strong sticky tape \textbullet{} a sunny window (and a dark cupboard for the class's second bag)

\begin{center}
\begin{tikzpicture}[scale=0.78]
  \draw[eq] (0,0) -- (0,6);
  \draw[Ink, line width=0.4pt] (-0.25,0) -- (-0.25,6);
  \draw[eq] (0.15,0.6) rectangle (3.4,4.6);
  \draw[Ink, line width=1pt] (0.15,4.3) -- (3.4,4.3);
  \fill[black!12] (0.2,0.65) rectangle (3.35,1.4);
  \draw[Ink, line width=0.8pt] (0.2,1.4) -- (3.35,1.4);
  \draw[Ink, line width=1.6pt] (0.6,4.75) -- (1.4,4.45) (2.2,4.45) -- (3.0,4.75);
  \draw[eq] (2.6,3.4) circle (0.3);
  \draw[eq] (0.8,3.3) .. controls (0.8,3.7) and (1.6,3.7) .. (1.6,3.3) -- cycle;
  \foreach \p in {(0.6,3.95),(1.2,4.05),(1.9,3.95),(2.9,4.0),(3.15,2.6)}{\fill[Ink] \p circle (0.05);}
  \eqlab{(3.25,4.6)}{(4.6,5.2)}{zip seal}
  \eqlab{(2.9,4.0)}{(4.6,4.3)}{droplets}
  \eqlab{(2.85,3.45)}{(4.6,3.5)}{sun and cloud drawn in marker}
  \eqlab{(3.35,1.0)}{(4.6,1.0)}{100\,mL coloured water}
  \eqlab{(1.0,4.6)}{(1.4,5.6)}{sticky tape}
  \eqlabl{(-0.25,2.5)}{(-0.9,2.5)}{window}
\end{tikzpicture}
\expfigcap{Scientific diagram of the water cycle bag taped to a window.}
\end{center}

\begin{infobox}{Safety}
Seal the bag fully and tape it firmly so it cannot fall. Food colouring stains: wipe up drips straight away.
\end{infobox}

\subsubsection*{Method}
\begin{enumerate}[itemsep=2pt]
  \item With the marker, draw a sun and a cloud near the top of the bag, and waves near the bottom.
  \item Measure 100\,mL of water, add two drops of food colouring, and pour it into the bag.
  \item Squeeze out most of the air and seal the bag. Check it does not leak.
  \item Tape the bag to a sunny window, seal at the top. Your teacher puts a second bag in a dark cupboard.
  \item Observe both bags straight away and at the start of each Science lesson for a week. Record each observation in the table.
\end{enumerate}

% ─── RESULTS ─────────────────────────────────────────────────────────────────
\exppart{Results}{Amber}

\subsubsection*{Observations of the two bags \skr{PC6}}
\begin{ctab}
\begin{tabular}{|SW{2.6cm}|W{2.6cm}|W{5.0cm}|W{4.8cm}|}
  \hline\hd & \thead{Day and time} & \thead{Weather} & \thead{Bag on the sunny window} & \thead{Bag in the cupboard} \\ \hline
  & \scaf{e.g. Day 0, 9:10 am} & \scaf{Sunny, warm} & \scaf{No droplets; water clear blue} & \scaf{No droplets} \\ \hline
  & Day 0 & & & \\ \hline
  & Day 1 & & & \\ \hline
  & Day 2 & & & \\ \hline
  & Day 3 & & & \\ \hline
  & Day 4 & & & \\ \hline
\end{tabular}
\end{ctab}

\needspace{14\baselineskip}
\task Sketch the sunny bag on the first day and on the last day. Label the droplets, the water and any drops running down. \skr{CO2}

\noindent
\begin{minipage}[t]{0.48\linewidth}\drawbox{6cm}{}\nameline{Day 0}\end{minipage}\hfill
\begin{minipage}[t]{0.48\linewidth}\drawbox{6cm}{}\nameline{Last day}\end{minipage}

% ─── ANALYSIS ────────────────────────────────────────────────────────────────
\exppart{Analysis and evaluation}{Coral}

\task Describe where in the sunny bag you saw evidence of evaporation, condensation and precipitation. Quote your observations. \skr{PA8}
\answerbox{5cm}

\needspace{10\baselineskip}
\task Compare the sunny bag with the bag in the cupboard. Explain the difference, using the word \emph{energy}. \skr{EV2}
\answerbox{5cm}

\needspace{12\baselineskip}
\task Look at the droplets at the top of the sunny bag.
\pt{a}Were they blue, or clear?
\answerbox{1cm}
\pt{b}Explain why. Use this to explain why rain is fresh water even when it comes from the salty sea.
\answerbox{5cm}

\needspace{10\baselineskip}
\task The amount of water at the bottom of the bag hardly changed all week. Explain why, using the word \emph{cycle}.
\answerbox{5cm}

\needspace{14\baselineskip}
\task Evaluate the bag as a model of the water cycle. \skr{EV3}
\pt{a}Give one strength of the model.
\answerbox{5cm}
\pt{b}Give two limitations of the model.
\answerbox{5cm}

\begin{yr8box}
\begin{keepext}
\exttask Plants release water vapour from their leaves (transpiration). Design a change to the bag model that would show transpiration as well, and predict what you would observe.
\answerbox{5cm}
\end{keepext}
\begin{keepext}
\exttask Earth's water cycle has been recycling the same water for billions of years. Use the bag model to explain how this is possible.
\answerbox{5cm}
\end{keepext}
\end{yr8box}

% ─── SKILL CHECK ─────────────────────────────────────────────────────────────
\exppart{Skill check}{NavyBlue}

\begin{skillcheck}
\scrow{QP1}{I can predict and explain what will happen.}
\scrow{PC6}{I recorded precise observations at each lesson, with the date and weather.}
\scrow{PA4}{I built a working model of the water cycle.}
\scrow{PA8}{I can match each part of the model to a stage of the water cycle.}
\scrow{EV2}{I can compare the two bags and explain the difference.}
\scrow{EV3}{I can give a strength and a limitation of the model.}
\scrow{CO2}{I can draw and label the set-up as a scientific diagram.}
\end{skillcheck}

\task Which skill will you focus on in Experiment 2? Explain what you will do differently.
\answerbox{5cm}
% ══════════════════════════════════════════════════════════════════════════════
%  EXPERIMENT 2 — HOUSEHOLD ENERGY AUDIT                                (Lesson 6)
%  Classroom survey in the lesson; the home audit is homework.
% ══════════════════════════════════════════════════════════════════════════════
\experimentdiv{Household Energy Audit}{Power ratings \textbullet{} Energy in kWh \textbullet{} Season and building design \textbullet{} A recommendation}{6}\label{sec:2}
\lessoncodes{VC2S8U16 \textbullet{} VC2S8H04 \textbullet{} VC2S8I02, I04, I05, I07, I08}

\vocabhead
\begin{ctab}
\begin{tabular}{|SW{4.2cm}|W{11.6cm}|}
  \hline\hd & \thead{Word} & \thead{Definition (copy from the slide)} \\ \hline
  & \scaf{\textbf{Energy audit}\label{voc:energy-audit}} & \scaf{A survey of where energy is used, how much, and how it could be reduced.} \\ \hline
  & \textbf{Power}\label{voc:power} & \\ \hline
  & \textbf{Kilowatt-hour}\label{voc:kilowatt-hour} & \\ \hline
  & \textbf{Energy rating label}\label{voc:energy-rating-label} & \\ \hline
  & \textbf{Standby}\label{voc:standby} & \\ \hline
\end{tabular}
\end{ctab}

\begin{lipanel}
\inquiryli{Collect data in a table, find the biggest energy users, and make a recommendation.}
\understandli{Explain how energy use depends on appliances, building design, season and climate.}
\sctitle
\begin{itemize}[leftmargin=1.4em, itemsep=0pt, topsep=1pt]
  \item I can calculate the energy an appliance uses in a day, in kWh.
  \item I can identify the biggest energy users in a home and explain why.
  \item I can write a recommendation to a household, backed by my data.
\end{itemize}
\end{lipanel}

\begin{skillsfocus}
\focus{PC5}{Audit a home with permission and privacy.}
\focus{PA1}{Organise appliance data in a table.}
\focus{PA2}{Draw a column graph of energy use.}
\focus{PA3}{Calculate energy $=$ power $\times$ time.}
\focus{PA6}{Does colder weather cause more use?}
\focus{EV4}{Judge energy advice and labels.}
\focus{CO1}{Write a recommendation backed by data.}
\end{skillsfocus}

% ─── TECHNIQUE ────────────────────────────────────────────────────────────────
\exppart{Technique spotlight}{Blue}

\begin{technique}{Reading power and energy labels}{PC6}
\begin{minipage}[c]{0.6\linewidth}
\textbf{Power} is how fast an appliance uses energy. It is printed on the appliance's rating plate, in watts (W) or kilowatts (kW): 1\,kW $=$ 1000\,W.

\smallskip
Many appliances also carry an \textbf{energy rating label}. More stars means the appliance uses less energy than others of its size. The number on the label is its estimated energy use in a year, in kilowatt-hours (kWh).
\end{minipage}\hfill
\begin{minipage}[c]{0.36\linewidth}\centering
\begin{tikzpicture}[scale=0.85]
  \draw[eq, rounded corners=2mm] (0,0) rectangle (4,4.2);
  \node[font=\headingfont\bfseries\footnotesize] at (2,3.85) {ENERGY RATING};
  \foreach \k in {0,...,5}{
    \pgfmathsetmacro\ang{160-\k*28}
    \draw[Ink, line width=0.8pt, fill=\ifnum\k<4 Ink\else white\fi] ($(2,1.6)+(\ang:1.6)+(90:0.2)$) -- ($(2,1.6)+(\ang:1.6)+(162:0.2)$) -- ($(2,1.6)+(\ang:1.6)+(234:0.2)$) -- ($(2,1.6)+(\ang:1.6)+(306:0.2)$) -- ($(2,1.6)+(\ang:1.6)+(18:0.2)$) -- cycle;}
  \node[font=\headingfont\bfseries\Large] at (2,1.25) {350};
  \node[font=\scriptsize] at (2,0.75) {kWh per year};
  \node[font=\scriptsize] at (2,0.35) {4 stars out of 6};
\end{tikzpicture}
\end{minipage}
\end{technique}

\begin{technique}{Calculating the energy an appliance uses}{PA3}
\textbf{Energy (kWh)} $=$ \textbf{power (kW)} $\times$ \textbf{time (hours)}. \quad Change watts to kilowatts by dividing by 1000, and minutes to hours by dividing by 60.
\end{technique}

\begin{workedex}{A kettle}
A 2000\,W kettle is used for 15 minutes a day. How much energy does it use in a day?

\smallskip
\textbf{Power:} $2000 \div 1000 = 2$\,kW. \quad \textbf{Time:} $15 \div 60 = 0.25$\,h.\\
\textbf{Energy:} $2 \times 0.25 = $ \textbf{0.5\,kWh} per day.
\end{workedex}

\begin{wedo}{A television}
A 120\,W television is on for 4 hours a day. How much energy does it use in a day?
\workbox{3.5cm}
\end{wedo}

\begin{youdobox}[Energy per day]
\task Complete the table. Then find the cost of each appliance per day, if electricity costs 30 cents per kWh. \skr{PA3}
\begin{ctab}
\begin{tabular}{|SW{3.4cm}|Z{2.2cm}|Z{2.4cm}|Z{3.6cm}|Z{3.4cm}|}
  \hline\hd & \thead{Appliance} & \thead{Power (W)} & \thead{Time per day (h)} & \thead{Energy per day (kWh)} & \thead{Cost per day (cents)} \\ \hline
  & \scaf{e.g. kettle} & \scaf{2000} & \scaf{0.25} & \scaf{$2 \times 0.25 = 0.5$} & \scaf{$0.5 \times 30 = 15$} \\ \hline
  & LED light globe & 10 & 5 & & \\ \hline
  & Electric heater & 2000 & 3 & & \\ \hline
  & Phone charger & 5 & 3 & & \\ \hline
\end{tabular}
\end{ctab}
\end{youdobox}

\begin{technique}{Judging energy advice}{EV4}
Energy advice is everywhere: in advertisements, on websites, from friends. Before you trust it, ask: \textbf{Who} wrote it? \textbf{Why} --- are they selling something? \textbf{What evidence} does it give? \textbf{When} was it written? Government energy rating labels are tested in the same way for every appliance, so they let you compare fairly.
\end{technique}

\begin{workedex}{Two sources}
\emph{Source 1:} an advertisement says ``Our heater is the cheapest to run!'' \emph{Source 2:} a government website lists the yearly energy use of every heater, tested the same way.

\smallskip
Source 2 is more reliable: it has no product to sell, and every heater is tested by the same method, so the numbers can be compared.
\end{workedex}

\begin{wedo}{Two sources}
A friend says that leaving the TV on standby uses no electricity at all. A government energy website says standby can make up about 3\% of a home's electricity use. Which would you trust, and why?
\answerbox{4.5cm}
\end{wedo}

% ─── PLAN AND SAFETY / METHOD ────────────────────────────────────────────────
\exppart{Plan, safety and method}{Teal}

\textbf{Aim:} To find out where energy is used in our classroom and in a home, and to recommend ways to use less.

\begin{infobox}{Ethics and safety}
\textbf{Ethics:} only audit your own home, and only with permission from the adults who live there. Do not write down names, addresses or bills in this booklet. \skr{PC5}\par
\textbf{Safety:} only read labels you can see easily. Do not unplug, move or open any appliance, and never touch a switch or plug with wet hands. Ask an adult if a label is hidden.
\end{infobox}

\begin{infobox}{Typical power ratings --- use these if you cannot read a label}
\renewcommand{\arraystretch}{1.15}\small
\begin{tabular}{@{}l r @{\hspace{8mm}} l r @{\hspace{8mm}} l r@{}}
Fridge (average) & 150\,W & Oven & 2400\,W & LED globe & 10\,W \\
Television & 120\,W & Washing machine & 500\,W & Halogen globe & 50\,W \\
Kettle & 2000\,W & Clothes dryer & 2400\,W & Laptop & 40\,W \\
Toaster & 900\,W & Electric heater & 2000\,W & Desktop computer & 100\,W \\
Microwave & 1000\,W & Air conditioner & 2500\,W & Game console & 150\,W \\
\end{tabular}
\end{infobox}

\subsubsection*{Method --- Part 1: The classroom (in class)}
\begin{enumerate}[itemsep=2pt]
  \item List every appliance that uses electricity in the classroom, and how many there are.
  \item Find the power of each, from its label or the table of typical ratings.
  \item Estimate how many hours a day each is switched on. Record everything in Table 2.1.
\end{enumerate}

\subsubsection*{Method --- Part 2: Your home (homework)}
\begin{enumerate}[itemsep=2pt]
  \item With an adult, list ten appliances in your home and find the power of each.
  \item Ask how many hours a day each is used in winter and in summer. Record these in Table 2.2.
  \item Complete the building design checklist (Table 2.3) together.
  \item Back in class, calculate the energy each appliance uses per day.
\end{enumerate}

% ─── RESULTS ─────────────────────────────────────────────────────────────────
\exppart{Results}{Amber}

\subsubsection*{Table 2.1 --- The classroom \skr{PA1}}
\begin{ctab}
\begin{tabular}{|SW{4.4cm}|Z{1.6cm}|Z{2.6cm}|Z{3.0cm}|Z{3.4cm}|}
  \hline\hd & \thead{Appliance} & \thead{Number} & \thead{Power of each (W)} & \thead{Hours on per day} & \thead{Energy per day (kWh)} \\ \hline
  & \scaf{e.g. LED ceiling light} & \scaf{12} & \scaf{20} & \scaf{6} & \scaf{$12 \times 0.02 \times 6 = 1.44$} \\ \hline
  & & & & & \\ \hline
  & & & & & \\ \hline
  & & & & & \\ \hline
  & & & & & \\ \hline
\end{tabular}
\end{ctab}

\clearpage
\subsubsection*{Table 2.2 --- My home}
\begin{ctab}
\begin{tabular}{|SW{3.6cm}|Z{2.0cm}|Z{2.2cm}|Z{2.6cm}|Z{2.2cm}|Z{2.6cm}|}
  \hline\hd & \thead{Appliance} & \thead{Power (W)} & \thead{Winter hours per day} & \thead{Winter kWh per day} & \thead{Summer hours per day} & \thead{Summer kWh per day} \\ \hline
  & \scaf{e.g. heater} & \scaf{2000} & \scaf{4} & \scaf{8} & \scaf{0} & \scaf{0} \\ \hline
  & & & & & & \\ \hline
  & & & & & & \\ \hline
  & & & & & & \\ \hline
  & & & & & & \\ \hline
  & & & & & & \\ \hline
  & & & & & & \\ \hline
  & & & & & & \\ \hline
  & & & & & & \\ \hline
  & & & & & & \\ \hline
  & & & & & & \\ \hline
\end{tabular}
\end{ctab}

\clearpage
\subsubsection*{Table 2.3 --- Building design checklist}
\begin{ctab}
\begin{tabular}{|SW{6.0cm}|Z{2.2cm}|W{6.8cm}|}
  \hline\hd & \thead{Feature} & \thead{Yes / No / Not sure} & \thead{How it changes energy use} \\ \hline
  & \scaf{Curtains or blinds on the windows} & \scaf{Yes} & \scaf{Keep heat in on winter nights and out on summer days.} \\ \hline
  & Ceiling insulation & & \\ \hline
  & Draught seals on doors & & \\ \hline
  & Double-glazed windows & & \\ \hline
  & Main living area has north-facing windows & & \\ \hline
  & All light globes are LEDs & & \\ \hline
  & Solar panels on the roof & & \\ \hline
\end{tabular}
\end{ctab}

\needspace{16\baselineskip}
\task Draw a column graph of the winter energy use per day of your six biggest users from Table 2.2. Label both axes and give the graph a title. \skr{PA2}
\graphgrid{9cm}

% ─── ANALYSIS ────────────────────────────────────────────────────────────────
\exppart{Analysis and evaluation}{Coral}

\task Which three appliances in your home use the most energy in winter? Suggest why they use so much. \skr{PA5}
\answerbox{5cm}

\needspace{14\baselineskip}
\task Compare your home's energy use in winter and summer. \skr{PA6}
\pt{a}Total energy per day in winter and in summer:
\workbox{3cm}
\pt{b}Energy use goes up as the weather gets colder. Is the cold weather itself the cause? Explain what actually uses the energy.
\answerbox{5cm}

\needspace{10\baselineskip}
\task Melbourne has cold winters and hot summers. Explain how a home in Darwin, where it is hot all year, would use energy differently.
\answerbox{5cm}

\needspace{10\baselineskip}
\task Choose two features from Table 2.3. Explain how each one changes the energy a home uses for heating and cooling.
\answerbox{5cm}

\needspace{12\baselineskip}
\task A family is choosing between two fridges of the same size. Fridge A has 3 stars and uses 450\,kWh a year. Fridge B has 5 stars and uses 300\,kWh a year. \skr{EV4}
\pt{a}How much energy would Fridge B save in a year?
\workbox{2.5cm}
\pt{b}Explain how the energy rating label helps the family choose.
\answerbox{5cm}

\needspace{12\baselineskip}
\task Write a recommendation to your household: three changes that would reduce its energy use. Back each one with data from your audit, and write in the third person. \skr{CO1}

\scaf{The audit found that the household uses the most energy for \dots. It is recommended that \dots, because \dots}
\answerbox{8cm}

\needspace{10\baselineskip}
\task Write a new question about household energy use that you could investigate next. \skr{EV5}
\answerbox{3cm}

\begin{yr8box}
\begin{keepext}
\exttask A household replaces eight 50\,W halogen globes with 10\,W LED globes. The lights are on for 5 hours a day. Calculate the energy saved in a year, and the money saved at 30 cents per kWh.
\workbox{4cm}
\end{keepext}
\begin{keepext}
\exttask Solar panels make the most electricity in the middle of the day, but many homes use the most electricity in the evening. Suggest two ways a household could make better use of its solar panels.
\answerbox{5cm}
\end{keepext}
\end{yr8box}

% ─── SKILL CHECK ─────────────────────────────────────────────────────────────
\exppart{Skill check}{NavyBlue}

\begin{skillcheck}
\scrow{PC5}{I audited a home with permission, and kept private information out of my booklet.}
\scrow{PA1}{I organised my appliance data in tables with units in the headings.}
\scrow{PA2}{I drew a column graph with labelled axes and a title.}
\scrow{PA3}{I can calculate energy with power $\times$ time.}
\scrow{PA6}{I can tell whether one thing causes another, or only changes with it.}
\scrow{EV4}{I can judge whether energy advice can be trusted.}
\scrow{CO1}{I can write a recommendation backed by data.}
\end{skillcheck}

\task Fill in the \emph{End of unit} column of your skills tracker at the front of this booklet. Which skill has improved most in this unit, and which one will you work on next?
\answerbox{5cm}

\clearpage
\hypertarget{glossary}{}\pdfbookmark[1]{Glossary}{glossarybm}
\markboth{Glossary}{}{\headingfont\bfseries\LARGE Glossary}\label{sec:glossary}\par\vspace{1mm}\noindent\rule{\linewidth}{0.6pt}\par
Every key vocabulary word in this booklet, and the page where you wrote what it means.\par
\begin{ctab}\renewcommand{\arraystretch}{1.5}
\begin{tabular}{|L{6cm}|C{2cm}|}
\hline\thead{Word} & \thead{Page} \\ \hline
\glossrow{Condensation}{voc:condensation}
\glossrow{Energy audit}{voc:energy-audit}
\glossrow{Energy rating label}{voc:energy-rating-label}
\glossrow{Evaporation}{voc:evaporation}
\glossrow{Kilowatt-hour}{voc:kilowatt-hour}
\glossrow{Limitation}{voc:limitation}
\glossrow{Model}{voc:model}
\glossrow{Power}{voc:power}
\glossrow{Precipitation}{voc:precipitation}
\glossrow{Standby}{voc:standby}
\end{tabular}
\end{ctab}

\end{document}
